\documentclass[a4paper,12pt]{article}
\usepackage[utf8]{inputenc}
\usepackage[german]{babel}
\usepackage{amsmath}
\usepackage{amssymb}
\usepackage{amsthm}
\usepackage{geometry}
\usepackage{graphicx}
\usepackage[hidelinks]{hyperref}
\usepackage{listings}
\usepackage{xcolor}
\usepackage{enumitem}
\usepackage{microtype}
\usepackage{booktabs}
\usepackage{array}
\usepackage{tikz}
\usetikzlibrary{arrows.meta,positioning,calc}

\geometry{margin=2.5cm}

% ---------- Ãœberlaufschutz ----------
\sloppy
\emergencystretch=3em
\setlength{\tolerance}{2000}
\setlength{\hbadness}{10000}

\theoremstyle{definition}
\newtheorem{definition}{Definition}
\newtheorem{beispiel}{Beispiel}
\newtheorem{satz}{Satz}
\newtheorem{lemma}{Lemma}
\newtheorem{bemerkung}{Bemerkung}

\setlist[itemize,1]{label=--}

% ---------- Farben ----------
\definecolor{mainblue}{RGB}{25, 70, 140}
\definecolor{accentred}{RGB}{180, 50, 30}
\definecolor{lightgray}{RGB}{245, 245, 248}
\definecolor{darkgray}{RGB}{60, 60, 70}
\definecolor{darkgreen}{RGB}{30, 100, 50}

\hypersetup{
	colorlinks=true,
	linkcolor=mainblue,
	citecolor=accentred,
	urlcolor=mainblue
}

%  Code-Highlighting
\lstset{
	basicstyle=\ttfamily\small,
	keywordstyle=\color{blue},
	commentstyle=\color{gray},
	stringstyle=\color{red},
	numbers=left,
	numberstyle=\tiny,
	breaklines=true,
	breakatwhitespace=false,
	frame=single,
	columns=fullflexible,
	literate=%
	{Ö}{{\"O}}1
	{Ä}{{\"A}}1
	{Ü}{{\"U}}1
	{ß}{{\ss}}1
	{ü}{{\"u}}1
	{ä}{{\"a}}1
	{ö}{{\"o}}1
	{Σ}{{$\Sigma$}}1
	{σ}{{$\sigma$}}1
	{δ}{{$\delta$}}1
	{→}{{$\rightarrow$}}1
	{⊆}{{$\subseteq$}}1
	{⊇}{{$\supseteq$}}1
	{≥}{{$\geq$}}1
}

\title{Sequenzgebundene REST-Schnittstellen (SC-REST):\\
Ein formales Erweiterungsmuster für die Abbildung\\
von Geschäftsprozessen auf REST-APIs}
\author{Stephan Epp}
\date{15. Juni 2026}

\begin{document}

\maketitle

\tableofcontents
\newpage

\section{Einführung}

REST (\emph{Representational State Transfer}) ist nach Fielding~\cite{fielding2000} ein
Architekturstil, dessen zentrales Prinzip die \emph{Zustandslosigkeit} (\emph{Statelessness})
zwischen aufeinanderfolgenden Anfragen ist: Jede Anfrage muss alle zu ihrer Bearbeitung
notwendigen Informationen enthalten, der Server hält keinen Sitzungszustand zwischen
Requests vor~\cite{rfc7231}. Dieses Prinzip ist wesentlich für Skalierbarkeit, Cachebarkeit
und die unabhängige Entwicklung von Client und Server.

In der praktischen Anwendung treten jedoch häufig \emph{Geschäftsprozesse} auf, die aus
fachlicher Sicht eine \emph{verbindliche Reihenfolge} von API-Aufrufen erfordern: Ein
Warenkorb muss erstellt werden, bevor Artikel hinzugefügt werden können; eine Zahlung kann
erst nach Abschluss des Checkout-Vorgangs ausgelöst werden; eine Bestellung kann erst
bestätigt werden, nachdem die Zahlung autorisiert wurde. Die klassische REST-Architektur
trifft hierzu keine normative Aussage: Ein zustandsloser Server kann grundsätzlich nicht
unterscheiden, ob eine Anfrage „zur falschen Zeit“ eintrifft, sofern die Anfrage selbst
syntaktisch und semantisch valide ist.

\subsection{Problemstellung}

Gegeben sei ein Geschäftsprozess, der durch eine Menge von API-Operationen
$\Sigma = \{\sigma_1, \ldots, \sigma_m\}$ (HTTP-Methode und Ressourcenpfad) realisiert wird.
Für eine konkrete \emph{Prozessinstanz} (z.\,B. eine Bestellung, einen Antrag, einen
Vertragsabschluss) existiert eine Menge \emph{zulässiger Aufrufsequenzen}
$c = (\sigma_{i_1}, \sigma_{i_2}, \ldots, \sigma_{i_k}) \in \Sigma^*$.

\textbf{Ziel:} Definiere eine formale, REST-kompatible Erweiterung, mit der ein Server
\begin{itemize}
	\item jede eingehende Anfrage gegen die für die jeweilige Prozessinstanz aktuell
	      zulässigen Operationen prüft,
	\item bei Abweichung von der vereinbarten Sequenz eine eindeutige Fehlermeldung liefert,
	      ohne die Zustandslosigkeit auf Protokollebene (jede Anfrage bleibt vollständig
	      selbstbeschreibend) aufzugeben,
	\item dem Client mittels \emph{Hypermedia as the Engine of Application State}
	      (HATEOAS) jederzeit die Menge der zulässigen Folgeoperationen mitteilt,
	\item die Kenntnis der wahrscheinlichen Folgeoperation(en) zur \emph{spekulativen
	      Vorbereitung} (Prefetching, Vorberechnung, Caching) nutzt.
\end{itemize}

\subsection{Motivation}

Drei Beobachtungen motivieren diese Arbeit:

\begin{enumerate}
	\item \textbf{Führung des Clients.} Ein Client, der eine ungültige Operation
	      versucht, erhält sofort eine maschinenlesbare Fehlermeldung mit der Menge der
	      tatsächlich zulässigen Folgeoperationen. Dies reduziert Implementierungsfehler
	      und Dokumentationsaufwand, da der Prozess \emph{im Protokoll selbst} kodiert ist.
	\item \textbf{Führung des Servers.} Da der Server zu jedem Zeitpunkt $T_{\text{think}}$
	      des Clients die Menge der \emph{möglichen} Folgezustände $N(s)$ kennt, kann er
	      Ressourcen, Datenbankabfragen oder Berechnungen für diese Zustände bereits während
	      der Bedenkzeit des Clients vorbereiten.
	\item \textbf{Strukturelle Nähe zum Subgraph-Problem.} Die Prüfung, ob eine beobachtete
	      Aufrufsequenz mit einem vorgegebenen Prozessautomaten konform ist, lässt sich als
	      Spezialfall eines Graph-Enthaltenseins-Problems auffassen. Für Prozessphasen mit
	      \emph{permutierbaren} optionalen Teilschritten erweist sich der in
	      \cite{epp2026subgraph} entwickelte Subgraph Algorithmus mit zyklischer Rotation
	      als das strukturell passende und nachweislich optimale Werkzeug.
\end{enumerate}

\subsection{Beitrag dieser Arbeit}

Diese Arbeit liefert:
\begin{itemize}
	\item eine formale Definition eines \emph{Prozessautomaten} als Erweiterung des
	      REST-Architekturstils (Abschnitt~\ref{sec:grundlagen}),
	\item die Definition des \emph{SC-REST-Patterns} (\emph{Sequence-Constrained REST})
	      mit konkreter HTTP-Semantik, einschließlich eines neuen Problem-Detail-Typs
	      gemäß RFC~7807~\cite{rfc7807} (Abschnitt~\ref{sec:pattern}),
	\item den formalen Nachweis von Korrektheit (Soundness/Completeness),
	      Laufzeitkomplexität und Optimalität der Sequenzprüfung
	      (Abschnitt~\ref{sec:analyse}),
	\item eine formale Anbindung permutierbarer optionaler Teilschritte an den Subgraph
	      Algorithmus aus~\cite{epp2026subgraph}, einschließlich eines neuen
	      Zustandsraum-Explosions-Lemmas (Abschnitt~\ref{sec:permutierbar}),
	\item ein formales Latenzmodell für serverseitiges Prefetching mit Beweis des
	      Latenzgewinns (Abschnitt~\ref{sec:prefetch}),
	\item eine Referenzimplementierung in Python (Abschnitt~\ref{sec:implementierung}),
	\item eine Veranschaulichung aller Ergebnisse durch matplotlib-Diagramme,
	\item einen ausführlichen Ausblick auf Standardisierung, verteilte Prozesse,
	      Verifikation und adaptive Vorhersagemodelle (Abschnitt~\ref{sec:ausblick}).
\end{itemize}

\section{Grundlagen}
\label{sec:grundlagen}

\subsection{REST und Hypermedia}

REST-APIs adressieren \emph{Ressourcen} über URIs und manipulieren diese mittels einer
einheitlichen Schnittstelle aus HTTP-Methoden (GET, POST, PUT, PATCH, DELETE). Im
\emph{Richardson Maturity Model}~\cite{fowler2010} bildet \emph{HATEOAS} die höchste Stufe:
Antworten enthalten Links auf zulässige Folgeoperationen, sodass der Client nicht auf
statische, externe Dokumentation angewiesen ist. SC-REST baut unmittelbar auf dieser Stufe
auf, erweitert sie jedoch um eine \emph{normative, serverseitig durchgesetzte} Komponente:
Die in den Links angegebenen Folgeoperationen sind nicht nur \emph{Vorschläge}, sondern die
\emph{einzigen} Operationen, die der Server für die betreffende Prozessinstanz akzeptiert.

\subsection{Formale Definitionen}

\begin{definition}[Operation]
	Eine \emph{Operation} $\sigma = (M, P)$ besteht aus einer HTTP-Methode
	$M \in \{\text{GET}, \text{POST}, \text{PUT}, \text{PATCH}, \text{DELETE}\}$ und einem
	Ressourcenpfad-Muster $P$. Die Menge aller im Prozess verwendeten Operationen wird mit
	$\Sigma$ bezeichnet (das \emph{Operationsalphabet}).
\end{definition}

\begin{definition}[Prozessautomat]
	\label{def:prozessautomat}
	Ein \emph{Prozessautomat} ist ein Tupel
	\[
		P = (S, \Sigma, \delta, s_0, F)
	\]
	mit:
	\begin{itemize}
		\item $S$: endliche Menge von \emph{Prozesszuständen},
		\item $\Sigma$: endliches Operationsalphabet,
		\item $\delta: S \times \Sigma \rightharpoonup S$: \emph{partielle} (deterministische)
		      Ãœbergangsfunktion,
		\item $s_0 \in S$: Startzustand,
		\item $F \subseteq S$: Menge der \emph{Endzustände} (abgeschlossene Prozessinstanzen).
	\end{itemize}
	Für $s \in S$ heißt
	\[
		N(s) := \{\, \delta(s,\sigma) \mid \sigma \in \Sigma,\ \delta(s,\sigma) \text{ ist definiert} \,\}
	\]
	die Menge der \emph{erreichbaren Folgezustände} und
	\[
		\Sigma(s) := \{\, \sigma \in \Sigma \mid \delta(s,\sigma) \text{ ist definiert} \,\}
	\]
	die Menge der in $s$ \emph{zulässigen Operationen}. Die Größe $d(s) := |\Sigma(s)|$
	heißt \emph{Verzweigungsgrad} von $s$.
\end{definition}

\begin{definition}[Aufrufsequenz und Lauf]
	\label{def:lauf}
	Eine \emph{Aufrufsequenz} $c = (c_1, \ldots, c_k) \in \Sigma^*$ ist eine endliche Folge
	von Operationen. Ein \emph{Lauf} von $P$ auf $c$ ist eine Folge von Zuständen
	$s_0, s_1, \ldots, s_k$ mit $s_i = \delta(s_{i-1}, c_i)$ für $i = 1, \ldots, k$.
\end{definition}

\begin{definition}[Konformität]
	\label{def:konformitaet}
	Eine Aufrufsequenz $c = (c_1, \ldots, c_k)$ heißt \emph{(partiell) konform} bezüglich
	$P$, wenn ein Lauf $s_0, \ldots, s_k$ gemäß Definition~\ref{def:lauf} existiert, d.\,h.
	wenn $\delta(s_{i-1}, c_i)$ für alle $i = 1, \ldots, k$ definiert ist. Die Sequenz heißt
	\emph{vollständig konform} (oder \emph{abgeschlossen}), wenn zusätzlich $s_k \in F$ gilt.
	Die Menge aller partiell konformen Sequenzen heißt $\mathcal{L}^*(P)$, die Menge aller
	vollständig konformen Sequenzen $\mathcal{L}(P) \subseteq \mathcal{L}^*(P)$.
\end{definition}

\begin{bemerkung}
	Da $\delta$ als \emph{Funktion} (nicht als Relation) definiert ist, ist der Folgezustand
	$s_i = \delta(s_{i-1}, c_i)$ für jeden Schritt eindeutig bestimmt. Diese Eigenschaft ist
	notwendig (aber nicht hinreichend), damit ein Server zu jedem Zeitpunkt einen
	wohldefinierten Prozesszustand pro Prozessinstanz führen kann.
\end{bemerkung}

\begin{beispiel}[Bestellprozess]
	\label{bsp:bestellprozess}
	Abbildung~\ref{fig:automat} zeigt einen Prozessautomaten für einen vereinfachten
	Bestellprozess mit
	$S = \{s_0, s_1, s_2, s_3, s_4, s_5\}$, $F = \{s_5\}$ und
	$\Sigma = \{$\texttt{POST /cart}, \texttt{POST /cart/items}, \texttt{POST /checkout},
	\texttt{POST /payment}, \texttt{POST /confirm}$\}$. Tabelle~\ref{tab:delta} listet die
	Übergangsfunktion $\delta$ vollständig auf. Jede Operation, die in einem Zustand $s$
	nicht in $\Sigma(s)$ enthalten ist, führt — gemäß der in
	Abschnitt~\ref{sec:pattern} definierten Semantik — zu einer
	\texttt{409 Conflict}-Antwort, ohne dass sich der Prozesszustand ändert.

	\begin{figure}[h]
		\centering
		\includegraphics[width=0.95\textwidth]{plots/plot1_automat.pdf}
		\caption{Prozessautomat eines Bestellprozesses. Durchgezogene blaue Pfeile
		bilden den (eindeutigen) Lauf des Hauptpfades; die Selbstschleife an $s_2$
		erlaubt das mehrfache Hinzufügen von Artikeln. Gestrichelte rote Pfeile
		repräsentieren Operationen, die in dem jeweiligen Zustand \emph{nicht} in
		$\Sigma(s)$ enthalten sind und daher gemäß Definition~\ref{def:sc-rest-server}
		mit \texttt{409 Conflict} beantwortet werden, ohne den Zustand zu verändern
		(der Zustand $\bot$ ist daher kein Element von $S$, sondern eine grafische
		Sammeldarstellung aller abgewiesenen Übergänge).}
		\label{fig:automat}
	\end{figure}

	\begin{table}[h]
		\centering
		\caption{Ãœbergangsfunktion $\delta$ des Bestellprozesses aus
		Abbildung~\ref{fig:automat}.}
		\label{tab:delta}
		\begin{tabular}{@{}lll@{}}
			\toprule
			Zustand $s$ & Operation $\sigma$ & Folgezustand $\delta(s,\sigma)$ \\
			\midrule
			$s_0$ & \texttt{POST /cart}       & $s_1$ \\
			$s_1$ & \texttt{POST /cart/items} & $s_2$ \\
			$s_2$ & \texttt{POST /cart/items} & $s_2$ \\
			$s_2$ & \texttt{POST /checkout}   & $s_3$ \\
			$s_3$ & \texttt{POST /payment}    & $s_4$ \\
			$s_4$ & \texttt{POST /confirm}    & $s_5$ \\
			\bottomrule
		\end{tabular}
	\end{table}

	Die Sequenz
	\[
		\begin{aligned}
			c = (&\texttt{POST /cart},\ \texttt{POST /cart/items},\ \texttt{POST /checkout},\\
			     &\texttt{POST /payment},\ \texttt{POST /confirm})
		\end{aligned}
	\]
	ist vollständig konform: Der zugehörige Lauf ist $s_0, s_1, s_2, s_3, s_4, s_5$ mit
	$s_5 \in F$. Die Sequenz $c' = (\texttt{POST /cart}, \texttt{POST /checkout})$ ist
	\emph{nicht} konform, da $\delta(s_1, \texttt{POST /checkout})$ undefiniert ist
	($\Sigma(s_1) = \{\texttt{POST /cart/items}\}$).
\end{beispiel}

\section{Das SC-REST-Pattern}
\label{sec:pattern}

\subsection{Prozessinstanzen und Serverzustand}

REST verlangt, dass \emph{jede Anfrage} alle notwendigen Informationen enthält. SC-REST
verletzt dieses Prinzip nicht: Der pro Prozessinstanz geführte Zustand $s \in S$ ist kein
verstecktes Sitzungsobjekt, sondern wird als \emph{Ressourcenattribut} der Prozessinstanz
modelliert und ist über deren URI adressierbar.

\begin{definition}[Prozessinstanz]
	\label{def:prozessinstanz}
	Eine \emph{Prozessinstanz} ist ein Tupel $(id, s)$, wobei $id$ ein eindeutiger
	Bezeichner (z.\,B. eine UUID, kodiert in der Ressourcen-URI
	\texttt{/orders/\{id\}}) und $s \in S$ der aktuelle Zustand der Instanz im
	Prozessautomaten $P$ ist. Der Server führt eine (partielle) Abbildung
	$\pi: \text{ID} \rightharpoonup S$ mit $\pi(id) = s$.
\end{definition}

\begin{definition}[SC-REST-Operationssemantik]
	\label{def:sc-rest-server}
	Sei $P = (S,\Sigma,\delta,s_0,F)$ ein Prozessautomat und $\pi(id) = s$ der aktuelle
	Zustand einer Prozessinstanz. Auf eine eingehende Anfrage $\sigma \in \Sigma$
	reagiert ein SC-REST-Server wie folgt:
	\begin{enumerate}
		\item \textbf{Falls $\delta(s,\sigma)$ definiert ist} (\emph{konformer Ãœbergang}):
		\begin{itemize}
			\item Führe die durch $\sigma$ beschriebene Fachlogik aus,
			\item setze $\pi(id) := \delta(s,\sigma)$,
			\item antworte mit einem 2xx-Statuscode und einer Repräsentation, die
			      gemäß HATEOAS Links auf alle Operationen
			      $\Sigma(\delta(s,\sigma))$ enthält.
		\end{itemize}
		\item \textbf{Falls $\delta(s,\sigma)$ undefiniert ist} (\emph{Sequenzverletzung}):
		\begin{itemize}
			\item $\pi(id)$ bleibt unverändert,
			\item der Server antwortet mit \texttt{409 Conflict} und einem
			      Problem-Detail-Objekt gemäß RFC~7807~\cite{rfc7807} vom Typ
			      \texttt{"{}sequence-violation"{}}, das die Menge $\Sigma(s)$ der
			      tatsächlich zulässigen Operationen enthält.
		\end{itemize}
	\end{enumerate}
\end{definition}

\begin{lstlisting}[language=, caption=Problem-Detail-Antwort bei Sequenzverletzung (RFC 7807)]
HTTP/1.1 409 Conflict
Content-Type: application/problem+json

{
  "type": "https://api.example.com/probs/sequence-violation",
  "title": "Ungueltige Operation fuer aktuellen Prozesszustand",
  "status": 409,
  "instance": "/orders/7f3e-...-91ab",
  "currentState": "s1",
  "attemptedOperation": "POST /checkout",
  "allowedOperations": [
    { "method": "POST", "path": "/orders/7f3e-.../items",
      "rel": "add-item" }
  ]
}
\end{lstlisting}

\begin{lstlisting}[language=, caption=Erfolgsantwort mit HATEOAS-Links auf zulaessige Folgeoperationen]
HTTP/1.1 201 Created
Content-Type: application/json

{
  "orderId": "7f3e-...-91ab",
  "state": "s2",
  "_links": {
    "add-item":  { "method": "POST",
                    "href": "/orders/7f3e-.../items" },
    "checkout":  { "method": "POST",
                    "href": "/orders/7f3e-.../checkout" }
  }
}
\end{lstlisting}

\subsection{Zustandslosigkeit und SC-REST}

\begin{bemerkung}
	SC-REST hebt die Zustandslosigkeit der einzelnen Anfrage nicht auf. Der Zustand $s$
	ist eine \emph{persistente Eigenschaft der adressierten Ressource} (der
	Prozessinstanz), nicht ein an die Verbindung oder Sitzung gebundenes Artefakt. Jede
	Anfrage bleibt selbstbeschreibend: Methode, Pfad und Repräsentation sind hinreichend,
	um die Anfrage zu bearbeiten; lediglich die \emph{Akzeptanz} der Anfrage hängt
	zusätzlich vom persistenten, ressourcengebundenen Zustand $\pi(id)$ ab — analog zur
	bekannten Bedingung, dass ein \texttt{PUT} auf eine Ressource mit
	\texttt{If-Match}-Header vom aktuellen \texttt{ETag} der Ressource abhängt. SC-REST
	verallgemeinert dieses Prinzip von einzelnen Versionskonflikten auf
	\emph{Sequenzkonflikte}.
\end{bemerkung}

\subsection{Implementierung der Übergangsprüfung}

Algorithmus~\ref{lst:transition} zeigt die Kernprüfung als Python-Funktion. Die
Ãœbergangsfunktion $\delta$ wird als verschachteltes \texttt{dict} (Hashtabelle)
$\Sigma \times S \to S$ realisiert, sodass die Prüfung in konstanter erwarteter Zeit
erfolgt (siehe Satz~\ref{satz:laufzeit-pruefung}).

\begin{lstlisting}[language=Python, label=lst:transition, caption=SC-REST Uebergangspruefung als Middleware-Funktion]
def handle_request(process: ProcessAutomaton,
                    instance_id: str,
                    operation: tuple) -> Response:
    """
    Prueft eine eingehende Operation gegen den aktuellen
    Zustand einer Prozessinstanz und fuehrt bei Konformitaet
    den Zustandsuebergang durch.

    delta: Dict[Tuple[State, Operation], State]
    """
    current_state = process.pi[instance_id]

    next_state = process.delta.get((current_state, operation))

    if next_state is None:
        # Sequenzverletzung: Zustand bleibt unveraendert
        allowed = [op for (s, op) in process.delta
                   if s == current_state]
        return problem_detail(
            status=409,
            type="sequence-violation",
            current_state=current_state,
            attempted=operation,
            allowed_operations=allowed,
        )

    # Konformer Uebergang: Fachlogik ausfuehren, Zustand setzen
    result = process.execute(operation, instance_id)
    process.pi[instance_id] = next_state

    return success_response(
        result=result,
        state=next_state,
        links=hateoas_links(process, next_state, instance_id),
    )
\end{lstlisting}

\section{Formale Analyse}
\label{sec:analyse}

Dieser Abschnitt erbringt vier formale Nachweise: (1) Korrektheit der Sequenzkontrolle
(Soundness und Completeness), (2) Laufzeit- und Optimalitätsschranken der
Übergangsprüfung, (3) die Charakterisierung der Konformitätsprüfung als
Pfad-Subgraph-Problem sowie deren Erweiterung auf permutierbare Teilschritte mittels des
Subgraph Algorithmus aus~\cite{epp2026subgraph}, und (4) ein formales Latenzmodell für
serverseitiges Prefetching.

\subsection{Korrektheit der Sequenzkontrolle}

\begin{satz}[Korrektheit der Sequenzkontrolle]
	\label{satz:korrektheit}
	Sei $P=(S,\Sigma,\delta,s_0,F)$ ein Prozessautomat und sei
	$c=(c_1,\ldots,c_k)\in\Sigma^*$ eine Aufrufsequenz, die einem SC-REST-Server gemäß
	Definition~\ref{def:sc-rest-server} beginnend bei $\pi(id)=s_0$ Operation für
	Operation präsentiert wird. Dann gilt:
	\begin{enumerate}
		\item[(a)] (\emph{Soundness}) Der Server antwortet auf $c_1,\ldots,c_k$ genau dann
		      \emph{niemals} mit \texttt{409 Conflict}, wenn $c \in \mathcal{L}^*(P)$
		      partiell konform ist.
		\item[(b)] (\emph{Completeness}) Ist $c$ partiell konform, so stimmt der nach
		      Verarbeitung von $c$ erreichte Serverzustand $\pi(id)$ exakt mit dem
		      eindeutigen Lauf-Endzustand $s_k$ aus Definition~\ref{def:lauf} überein.
		\item[(c)] (\emph{Fehlerlokalisierung}) Ist $c$ \emph{nicht} konform, so existiert
		      ein eindeutiges minimales $i^\ast \in \{1,\ldots,k\}$, sodass
		      $c_1,\ldots,c_{i^\ast-1}$ konform ist, $\delta(s_{i^\ast-1},c_{i^\ast})$
		      undefiniert ist, und der Server bei genau diesem Schritt -- und nur bei
		      diesem -- mit \texttt{409 Conflict} antwortet, wobei $\pi(id)$ auf
		      $s_{i^\ast-1}$ verbleibt.
	\end{enumerate}
\end{satz}

\begin{proof}
	Wir führen vollständige Induktion über $i=0,\ldots,k$ und zeigen die Invariante
	\[
		I(i):\quad \pi(id) \text{ ist nach Verarbeitung von } c_1,\ldots,c_i
		\text{ wohldefiniert und gleich } s_i,
	\]
	solange $c_1,\ldots,c_i$ konform ist, andernfalls gleich $s_{i^\ast-1}$ für das
	kleinste $i^\ast \le i$, für das $\delta(s_{i^\ast-1},c_{i^\ast})$ undefiniert ist.

	\textbf{Induktionsanfang ($i=0$):} $\pi(id)=s_0$ nach Voraussetzung; $I(0)$ gilt
	trivial, da die leere Sequenz konform ist mit Lauf $s_0$.

	\textbf{Induktionsschritt ($i-1 \to i$):} Nach Induktionsannahme gilt $I(i-1)$.
	Zwei Fälle:
	\begin{itemize}
		\item \emph{Fall A: $c_1,\ldots,c_{i-1}$ ist konform und $\pi(id)=s_{i-1}$.}
		      Bei Eintreffen von $c_i$ prüft der Server gemäß
		      Definition~\ref{def:sc-rest-server}, ob $\delta(s_{i-1},c_i)$ definiert ist.
		      \begin{itemize}
			      \item Ist es definiert, setzt der Server $\pi(id):=\delta(s_{i-1},c_i)=s_i$
			            (Definition~\ref{def:lauf}) und antwortet nicht mit \texttt{409}.
			            Damit ist $c_1,\ldots,c_i$ konform und $I(i)$ gilt mit $\pi(id)=s_i$.
			      \item Ist es undefiniert, antwortet der Server mit \texttt{409 Conflict}
			            und $\pi(id)$ bleibt bei $s_{i-1}$. Damit ist $i^\ast = i$ das
			            gesuchte minimale Indexelement, und $I(i)$ gilt mit
			            $\pi(id)=s_{i^\ast-1}=s_{i-1}$.
		      \end{itemize}
		\item \emph{Fall B: $c_1,\ldots,c_{i-1}$ ist bereits nicht konform mit
		      Verletzungsindex $i^\ast \le i-1$ und $\pi(id)=s_{i^\ast-1}$.} Da $\delta$ in
		      Definition~\ref{def:sc-rest-server} bei einer Sequenzverletzung den Zustand
		      \emph{nicht} verändert, bleibt $\pi(id)=s_{i^\ast-1}$ auch nach Verarbeitung
		      von $c_i$ -- unabhängig davon, ob $\delta(s_{i^\ast-1},c_i)$ definiert wäre,
		      denn der Server prüft stets gegen den \emph{aktuellen}, nicht den
		      ursprünglich vorgesehenen Zustand. $I(i)$ gilt also mit demselben $i^\ast$.
	\end{itemize}

	Aus $I(k)$ folgen unmittelbar (a)--(c): Tritt während der gesamten Verarbeitung von
	$c_1,\ldots,c_k$ kein Fall „undefiniert“ auf, ist $c$ nach Definition~\ref{def:lauf}
	konform mit Lauf $s_0,\ldots,s_k$ und der Server antwortet nie mit \texttt{409}
	(Soundness, durch Kontraposition von Fall A in beiden Richtungen).
	Tritt der Fall „undefiniert“ erstmals bei Index $i^\ast$ auf, so ist dies nach
	Konstruktion eindeutig und minimal, der Server antwortet exakt einmal -- bei
	$i^\ast$ -- mit \texttt{409}, und $\pi(id)$ verbleibt für alle $i \ge i^\ast$ bei
	$s_{i^\ast-1}$ (Fehlerlokalisierung, (c)). Im konformen Fall liefert $I(k)$ direkt
	(b). \qedhere
\end{proof}

\begin{bemerkung}
	Satz~\ref{satz:korrektheit}(c) ist die formale Grundlage für die in der Einführung
	geforderte \emph{Führung des Clients}: Der Server lehnt nicht nur ab, sondern tut dies
	beim \emph{frühestmöglichen} Zeitpunkt einer Sequenzverletzung, und der Zustand der
	Prozessinstanz bleibt für \emph{beliebig viele} nachfolgende fehlerhafte Anfragen
	stabil bei $s_{i^\ast-1}$ -- ein Client kann also nach einem \texttt{409} beliebig oft
	erneut die laut \texttt{allowedOperations} korrekte Operation versuchen, ohne dass
	der Prozesszustand durch die fehlgeschlagenen Versuche beeinflusst wird
	(\emph{Stabilität unter Fehlversuchen}).
\end{bemerkung}

\subsection{Laufzeit und Optimalität der Übergangsprüfung}

\begin{satz}[Laufzeit der Übergangsprüfung]
	\label{satz:laufzeit-pruefung}
	Wird $\delta$ als Hashtabelle $H: S \times \Sigma \to S \cup \{\text{undef}\}$
	realisiert (Algorithmus~\ref{lst:transition}), so benötigt die Prüfung, ob
	$\delta(s,\sigma)$ definiert ist, sowie -- falls definiert -- die Bestimmung von
	$\delta(s,\sigma)$, im Erwartungswert $O(1)$ Zeit, unabhängig von $|S|$, $|\Sigma|$
	und der Länge $k$ der bisherigen Aufrufsequenz.
\end{satz}

\begin{proof}
	Ein Hashtabellen-Lookup $H[(s,\sigma)]$ erfordert die Berechnung eines Hashwerts über
	das Paar $(s,\sigma)$ -- beides Werte konstanter Beschreibungsgröße (Zustands-ID bzw.
	Methode+Pfad-Muster) -- sowie den Zugriff auf den entsprechenden Bucket. Beide
	Operationen sind unter Standardannahmen an die Hashfunktion (gleichverteilte Streuung,
	Lastfaktor $O(1)$) im Erwartungswert $O(1)$. Die Anzahl $|S|$ der Zustände und
	$|\Sigma|$ der Operationen beeinflusst lediglich die \emph{Größe} der Hashtabelle,
	nicht jedoch -- bei konstantem Lastfaktor -- die erwartete Zugriffszeit. Die Länge $k$
	der bisherigen Sequenz geht nicht in die Prüfung ein, da nur der aktuelle Zustand
	$\pi(id)=s_{k}$ betrachtet wird (Markov-Eigenschaft des Prozessautomaten,
	Definition~\ref{def:lauf}). \qedhere
\end{proof}

\begin{satz}[Optimalität der Gesamtprüfung einer Sequenz]
	\label{satz:optimalitaet-sequenz}
	Die Prüfung der Konformität einer vollständigen Aufrufsequenz $c=(c_1,\ldots,c_k)$
	gemäß Definition~\ref{def:konformitaet} benötigt $\Theta(k)$ Zeit und ist
	asymptotisch optimal: Kein korrekter Algorithmus kann diese Prüfung in $o(k)$ Zeit
	durchführen.
\end{satz}

\begin{proof}
	\textbf{Obere Schranke $O(k)$.} Nach Satz~\ref{satz:laufzeit-pruefung} kostet jeder
	der $k$ Ãœbergangsschritte $O(1)$, in Summe also $O(k)$.

	\textbf{Untere Schranke $\Omega(k)$.} Sei $\mathcal{A}$ ein beliebiger korrekter
	Algorithmus, der für eine Eingabesequenz $c=(c_1,\ldots,c_k)$ entscheidet, ob
	$c \in \mathcal{L}^*(P)$. Angenommen, $\mathcal{A}$ liest ein Element $c_j$ für ein
	$j\in\{1,\ldots,k\}$ nicht. Wir konstruieren zwei Eingaben $c$ und $c'$, die sich nur
	in $c_j$ unterscheiden:
	\[
		c = (c_1,\ldots,c_{j-1}, \sigma, c_{j+1},\ldots,c_k), \qquad
		c' = (c_1,\ldots,c_{j-1}, \sigma', c_{j+1},\ldots,c_k),
	\]
	wobei $\sigma, \sigma' \in \Sigma(s_{j-1})$ mit $\sigma \neq \sigma'$ und
	$\delta(s_{j-1},\sigma) \neq \delta(s_{j-1},\sigma')$ -- eine solche Wahl existiert,
	sobald $d(s_{j-1}) \ge 2$, was o.\,B.\,d.\,A. für mindestens ein $j$ gilt (sonst ist
	$P$ trivial, da jeder Zustand höchstens einen Folgezustand besitzt, und die Aussage
	wird leer). Für geeignete Fortsetzungen $c_{j+1},\ldots,c_k$ -- etwa eine konforme
	Fortsetzung ab $\delta(s_{j-1},\sigma)$, die ab $\delta(s_{j-1},\sigma')$ \emph{nicht}
	konform ist -- unterscheiden sich die korrekten Antworten für $c$ und $c'$. Da
	$\mathcal{A}$ den Eintrag $c_j$ nicht liest, verhält sich $\mathcal{A}$ auf $c$ und
	$c'$ identisch und liefert für mindestens eine der beiden Eingaben ein falsches
	Ergebnis -- Widerspruch zur Korrektheit. Folglich muss $\mathcal{A}$ jedes Element
	$c_1,\ldots,c_k$ lesen, also $T_{\mathcal{A}}(k) \ge \Omega(k)$. \qedhere
\end{proof}

\begin{bemerkung}
	Satz~\ref{satz:optimalitaet-sequenz} zeigt, dass die SC-REST-Konformitätsprüfung --
	anders als das allgemeine Subgraph-Problem aus~\cite{epp2026subgraph} mit seiner
	unteren Schranke $\Omega(n^{3-\varepsilon})$ -- \emph{linear} und damit
	vergleichsweise „billig“ ist. Der Grund liegt in der \emph{linearen Struktur} einer
	Aufrufsequenz: Sie ist ein einfacher Pfad, kein beliebiger Graph. Abschnitt
	\ref{sec:permutierbar} zeigt, dass diese Eigenschaft verloren geht, sobald
	\emph{permutierbare} Teilschritte zugelassen werden -- und genau dann wird der
	Subgraph Algorithmus mit zyklischer Rotation wieder relevant.
\end{bemerkung}

\subsection{Konformität als Pfad-Subgraph-Problem}

\begin{definition}[Prozessgraph]
	\label{def:prozessgraph}
	Der \emph{Prozessgraph} $G_P=(S,E)$ zu einem Prozessautomaten
	$P=(S,\Sigma,\delta,s_0,F)$ ist der gerichtete, kantenbeschriftete Graph mit
	\[
		E = \{\, (s,s',\sigma) \mid s,s' \in S,\ \sigma\in\Sigma,\ \delta(s,\sigma)=s' \,\}.
	\]
	Eine Aufrufsequenz $c=(c_1,\ldots,c_k)$ induziert den \emph{Sequenzgraphen}
	$G_c$: einen einfachen, gerichteten Pfad mit $k+1$ Knoten $u_0,\ldots,u_k$ und Kanten
	$(u_{i-1},u_i,c_i)$ für $i=1,\ldots,k$.
\end{definition}

\begin{satz}[Pfad-Charakterisierung der Konformität]
	\label{satz:pfadcharakterisierung}
	Eine Aufrufsequenz $c$ ist genau dann partiell konform bezüglich $P$
	(Definition~\ref{def:konformitaet}), wenn der Sequenzgraph $G_c$ unter der Abbildung
	$u_0 \mapsto s_0$ als beschrifteter Teilpfad in $G_P$ einbettbar ist, d.\,h. wenn eine
	Folge von Knoten $s_0=v_0,v_1,\ldots,v_k \in S$ existiert mit
	$(v_{i-1},v_i,c_i)\in E$ für alle $i=1,\ldots,k$.
\end{satz}

\begin{proof}
	„$\Rightarrow$“: Ist $c$ konform, existiert nach Definition~\ref{def:lauf} ein Lauf
	$s_0,\ldots,s_k$ mit $\delta(s_{i-1},c_i)=s_i$. Nach Definition~\ref{def:prozessgraph}
	ist dann $(s_{i-1},s_i,c_i)\in E$ für alle $i$, also bildet $v_i:=s_i$ die geforderte
	Pfadeinbettung.

	„$\Leftarrow$“: Existiert eine solche Knotenfolge $v_0=s_0,\ldots,v_k$ mit
	$(v_{i-1},v_i,c_i)\in E$, so gilt nach Definition~\ref{def:prozessgraph}
	$\delta(v_{i-1},c_i)=v_i$ für alle $i$. Da $\delta$ eine Funktion ist
	(Definition~\ref{def:prozessautomat}), ist diese Folge der \emph{eindeutige} Lauf von
	$P$ auf $c$ mit $s_i=v_i$, also ist $c$ konform. \qedhere
\end{proof}

\begin{bemerkung}
	\label{bem:einfacher-fall}
	Satz~\ref{satz:pfadcharakterisierung} reduziert die Konformitätsprüfung auf die
	Einbettbarkeit eines \emph{einfachen, beschrifteten Pfades} $G_c$ in $G_P$ -- ein
	Spezialfall des allgemeinen Subgraph-Enthaltenseins-Problems. Während das allgemeine
	Problem (zwei beliebige Graphen, Einbettung unter Knotenpermutation bzw. zyklischer
	Rotation) nach~\cite{epp2026subgraph} eine scharfe untere Schranke
	$\Omega(n^{3-\varepsilon})$ unter SETH besitzt, ist die Pfad-Einbettung \emph{ab dem
	festen Startknoten $s_0$} wegen der Determiniertheit von $\delta$ trivial in
	$\Theta(k)$ entscheidbar (Satz~\ref{satz:optimalitaet-sequenz}): An jedem Knoten $v_i$
	existiert höchstens \emph{eine} mit $c_{i+1}$ beschriftete ausgehende Kante, sodass
	keine Suche über mehrere Kandidaten und keine Rotation erforderlich ist. Der
	allgemeine, schwere Fall des Subgraph-Problems tritt erst dann wieder auf, wenn der
	Startknoten \emph{nicht fest} ist oder wenn -- wie in
	Abschnitt~\ref{sec:permutierbar} -- mehrere Kanten mit unterschiedlichen Beschriftungen
	in beliebiger Reihenfolge durchlaufen werden dürfen.
\end{bemerkung}

\subsection{Permutierbare optionale Teilschritte}
\label{sec:permutierbar}

Viele Geschäftsprozesse enthalten Phasen, in denen mehrere optionale Teilschritte
\emph{in beliebiger Reihenfolge}, jedoch jeweils höchstens einmal, ausgeführt werden
dürfen, bevor der Prozess fortgesetzt wird -- etwa das Hinzufügen mehrerer unabhängiger
Zusatzoptionen zu einer Bestellung (Geschenkverpackung, Expresslieferung, Gutscheincode,
Versicherung), deren Reihenfolge fachlich irrelevant ist.

\begin{definition}[Sammelzustand mit permutierbaren Teilschritten]
	\label{def:sammelzustand}
	Ein Zustand $s_c \in S$ heißt \emph{Sammelzustand mit $r$ permutierbaren optionalen
	Teilschritten}, wenn ein Operationssatz
	$\Sigma_{\mathrm{opt}} = \{\sigma_1,\ldots,\sigma_r\} \subseteq \Sigma$ sowie eine
	Abschlussoperation $\sigma_{\mathrm{done}} \in \Sigma$ existieren, sodass für jede
	Teilmenge $T \subseteq \{1,\ldots,r\}$ der bereits ausgeführten Operationen und jedes
	$j \notin T$ gilt: Aus dem zu $T$ korrespondierenden (logischen) Zwischenzustand ist
	$\sigma_j$ zulässig und führt zum Zwischenzustand $T \cup \{j\}$; aus dem
	Zwischenzustand zu $T=\{1,\ldots,r\}$ (alle Teilschritte ausgeführt, oder ein
	beliebiges zulässiges $T$, falls $\sigma_{\mathrm{done}}$ auch früher erlaubt ist)
	führt $\sigma_{\mathrm{done}}$ in den Folgezustand des Prozesses.
\end{definition}

\begin{lemma}[Zustandsraum-Explosion bei naiver Automatenkodierung]
	\label{lem:zustandsexplosion}
	Sei $s_c$ ein Sammelzustand mit $r$ permutierbaren optionalen Teilschritten gemäß
	Definition~\ref{def:sammelzustand}, und seien für je zwei verschiedene Teilmengen
	$T \neq T' \subseteq \{1,\ldots,r\}$ die jeweils zulässigen Restoperationen
	$\Sigma_{\mathrm{opt}} \setminus T$ und $\Sigma_{\mathrm{opt}} \setminus T'$
	verschieden (d.\,h. $\Sigma(s_T) \neq \Sigma(s_{T'})$ für die zu $T$ bzw. $T'$
	gehörigen Zwischenzustände $s_T,s_{T'}$). Dann benötigt \emph{jeder} deterministische
	Prozessautomat, der diesen Sammelzustand abbildet, mindestens $2^r$ Zustände, also
	$|S| \geq 2^r$.
\end{lemma}

\begin{proof}
	Wir zeigen, dass die $2^r$ Zwischenzustände $s_T$, $T \subseteq \{1,\ldots,r\}$,
	paarweise \emph{verschieden} sein müssen -- ein Argument in der Tradition des
	Myhill-Nerode-Theorems für reguläre Sprachen, angewandt auf
	Erreichbarkeits-Äquivalenzklassen von Zuständen.

	Angenommen, $s_T = s_{T'}$ für zwei verschiedene Teilmengen $T \neq T'$. Da $\delta$
	als Funktion eines einzelnen Zustands definiert ist
	(Definition~\ref{def:prozessautomat}), folgt $\Sigma(s_T) = \Sigma(s_{T'})$ -- die
	Menge der von $s_T$ aus zulässigen Operationen hängt nur vom Zustand selbst ab, nicht
	von der „Historie“, über die er erreicht wurde. Dies widerspricht jedoch der
	Voraussetzung $\Sigma(s_T) \neq \Sigma(s_{T'})$ für $T \neq T'$.

	Also müssen alle $2^r$ Teilmengen $T \subseteq \{1,\ldots,r\}$ auf paarweise
	verschiedene Zustände $s_T \in S$ abgebildet werden, woraus $|S| \geq 2^r$ folgt.
	\qedhere
\end{proof}

\begin{bemerkung}
	Lemma~\ref{lem:zustandsexplosion} ist die formale Begründung dafür, dass ein
	SC-REST-Server für Sammelzustände \emph{nicht} naiv $\delta$ über den vollständigen
	Potenzmengen-Zustandsraum $2^{\{1,\ldots,r\}}$ kodieren sollte: Bereits für $r=20$
	permutierbare Teilschritte ergäbe dies über $10^6$ Zustände allein für diese eine
	Prozessphase. Stattdessen genügt es, gemäß Satz~\ref{satz:permutierbar-subgraph} eine
	kompakte Repräsentation der Größe $O(r)$ zu führen und die Konformität einer
	beobachteten Teilreihenfolge mittels des Subgraph Algorithmus zu prüfen.
\end{bemerkung}

\begin{definition}[Signaturfolgen für Sammelzustände]
	\label{def:signaturfolge-sammel}
	Für einen Sammelzustand $s_c$ mit Operationssatz
	$\Sigma_{\mathrm{opt}}=\{\sigma_1,\ldots,\sigma_r\}$ sei $\sigma^P =
	[\sigma_1,\ldots,\sigma_r]$ die \emph{kanonische Reihenfolge} (z.\,B. die im
	Prozessmodell dokumentierte Referenzreihenfolge). Eine vom Client beobachtete
	Ausführungsreihenfolge $c' = (c'_1,\ldots,c'_r)$ ist eine Permutation von
	$\Sigma_{\mathrm{opt}}$. Wir bezeichnen mit $\sigma^{C}=[c'_1,\ldots,c'_r]$ die
	zugehörige \emph{Beobachtungsfolge}.
\end{definition}

\begin{satz}[Verifikation permutierbarer Teilschritte via Subgraph Algorithmus]
	\label{satz:permutierbar-subgraph}
	Sei $s_c$ ein Sammelzustand mit $r$ permutierbaren optionalen Teilschritten und sei
	$c'=(c'_1,\ldots,c'_r)$ eine vom Client vollständig durchlaufene Beobachtungsfolge
	(jede Operation aus $\Sigma_{\mathrm{opt}}$ genau einmal, in beliebiger Reihenfolge).
	Dann gilt:
	\begin{enumerate}
		\item[(a)] Die Aussage „$c'$ ist eine zulässige Reihenfolge (Permutation) der
		      Operationen aus $\Sigma_{\mathrm{opt}}$“ ist äquivalent zur Aussage „die
		      Signaturfolge $\sigma^{C}$ ist -- als Multimenge -- identisch zu
		      $\sigma^{P}$, d.\,h. $\sigma^{C}$ entsteht aus $\sigma^{P}$ durch
		      Anwendung einer Permutation auf die Indexmenge $\{1,\ldots,r\}$“.
		\item[(b)] Beschränkt man die zulässigen Umordnungen auf \emph{zyklische
		      Rotationen} von $\sigma^{P}$ (z.\,B. weil eine Teilmenge der $r$
		      Teilschritte eine vom Prozessdesign vorgegebene relative Ordnung
		      beibehalten muss, die nur als Ganzes zyklisch verschoben werden darf), so
		      lässt sich die Prüfung von (a) mit dem Subgraph Algorithmus aus
		      \cite{epp2026subgraph} in $\Theta(r^3)$ Zeit durchführen, und diese
		      Laufzeit ist nach dem Optimalitätssatz von~\cite{epp2026subgraph}
		      asymptotisch optimal unter SETH: Keine Prüfung ist in
		      $O(r^{3-\varepsilon})$ Zeit für $\varepsilon>0$ möglich.
		\item[(c)] Im allgemeinen Fall (beliebige Permutationen statt zyklischer
		      Rotationen) genügt bereits eine Sortierung der Signaturfolgen in
		      $O(r \log r)$ Zeit, um (a) zu entscheiden.
	\end{enumerate}
\end{satz}

\begin{proof}
	\textbf{(a)} Da jede Operation aus $\Sigma_{\mathrm{opt}}$ in $c'$ nach Voraussetzung
	genau einmal vorkommt, ist $c'$ per Definition eine Permutation von
	$(\sigma_1,\ldots,\sigma_r)$. Die zugehörige Signaturfolge $\sigma^{C}$ entsteht somit
	aus $\sigma^{P}=[\sigma_1,\ldots,\sigma_r]$ durch Umordnung gemäß genau dieser
	Permutation. Umgekehrt definiert jede Umordnung von $\sigma^{P}$, die als Multimenge
	mit $\sigma^{P}$ übereinstimmt, eindeutig eine Permutation von
	$\{1,\ldots,r\}$ und damit eine zulässige Ausführungsreihenfolge im Sinne von
	Definition~\ref{def:sammelzustand} (jede Operation wird, unabhängig von der
	Reihenfolge, genau einmal ausgeführt und überführt den Sammelzustand in den
	Folgezustand des Prozesses).

	\textbf{(b)} Wende den Subgraph Algorithmus aus~\cite{epp2026subgraph} auf die beiden
	Sequenzen $\sigma^{P}$ und $\sigma^{C}$ an: Beide haben Länge $r$, die
	Spalten-Signaturen $\sigma_j = \sum_i A_{ij}2^i + j\cdot 2^r$ werden für beide Folgen
	gebildet (hier interpretiert als $1\times r$-„Adjazenzmatrizen“ über dem
	Operationsalphabet, wobei $A_{ij}=1$ gdw. die $i$-te Operation an Position $j$ der
	jeweiligen Folge steht). Für jede der $r$ zyklischen Rotationen von $\sigma^{C}$ wird
	via LCS geprüft, ob $\sigma^{P}\subseteq\mathrm{rot}_k(\sigma^{C})$ gilt
	(Satz~2 in~\cite{epp2026subgraph}). Nach Satz~„Laufzeit“ und
	Satz~„Optimalität des Subgraph Algorithmus“ aus~\cite{epp2026subgraph} ist diese
	Prüfung in $\Theta(r^3)$ durchführbar und unter SETH nicht in
	$O(r^{3-\varepsilon})$ für $\varepsilon>0$ verbesserbar -- die untere Schranke aus
	\cite{epp2026subgraph} (kombiniert aus Eingabe-Leseschranke, Notwendigkeit aller
	Rotationen und der SETH-basierten LCS-Schranke) übertragen sich unverändert, da die
	Signaturkonstruktion injektiv ist (Lemma „Eindeutigkeit der Signaturen“,
	\cite{epp2026subgraph}) und somit keine Information verloren geht.

	\textbf{(c)} Da Signaturen nach \cite{epp2026subgraph} (Lemma „Eindeutigkeit der
	Signaturen“) injektiv und damit als Schlüssel für eine Ordnungsrelation geeignet sind,
	kann man $\sigma^{C}$ und $\sigma^{P}$ jeweils in $O(r\log r)$ sortieren und
	elementweise vergleichen: Gleichheit der sortierten Folgen ist äquivalent zur
	Multimengen-Gleichheit aus (a). Dies ist günstiger als $\Theta(r^3)$, jedoch nur
	anwendbar, wenn \emph{beliebige} Permutationen zulässig sind und keine relative
	Ordnung zwischen Teilgruppen der $r$ Operationen erhalten bleiben muss. \qedhere
\end{proof}

\begin{bemerkung}
	\label{bem:permutierbar-fazit}
	Satz~\ref{satz:permutierbar-subgraph} zeigt eine doppelte Erweiterung gegenüber
	\cite{epp2026subgraph}: Erstens wird der Subgraph Algorithmus von seinem
	ursprünglichen Anwendungsfeld (AST-Vergleich, Graph-Deduplizierung) auf die
	Konformitätsprüfung permutierbarer API-Teilschritte \emph{übertragen}. Zweitens zeigt
	Lemma~\ref{lem:zustandsexplosion} zusammen mit Satz~\ref{satz:permutierbar-subgraph},
	dass die signaturbasierte Prüfung eine \emph{exponentielle} Raumersparnis erzielt: Statt
	$\Omega(2^r)$ Zuständen im naiv kodierten Automaten genügt eine Signaturfolge der
	Größe $O(r)$ zuzüglich einer Prüfung in $\Theta(r^3)$ -- ein klassischer
	Zeit-Raum-Trade-off, der für praktisch relevante $r$ (typischerweise $r \le 20$, siehe
	Abbildung~\ref{fig:komplexitaet}) deutlich zugunsten der signaturbasierten Methode
	ausfällt, da $2^{20} \approx 10^6 \gg 20^3 = 8000$.
\end{bemerkung}

\begin{figure}[h]
	\centering
	\includegraphics[width=\textwidth]{plots/plot3_komplexitaet.pdf}
	\caption{Links: Vergleich der Konformitätsprüfung einer festen, linearen
	Aufrufsequenz ($O(k)$, Satz~\ref{satz:optimalitaet-sequenz}) mit der signaturbasierten
	Prüfung permutierbarer Teilschritte mittels Subgraph Algorithmus und zyklischer
	Rotation ($\Theta(r^3)$, Satz~\ref{satz:permutierbar-subgraph}). Beide Kurven wachsen
	polynomial, jedoch mit unterschiedlichem Grad -- die lineare Sequenzprüfung bleibt
	für alle betrachteten Größen um Größenordnungen günstiger, was die in
	Bemerkung~\ref{bem:einfacher-fall} beschriebene strukturelle Vereinfachung
	widerspiegelt. Rechts: Vergleich des Subgraph Algorithmus ($\Theta(r^3)$) mit der
	naiven vollständigen Permutationsprüfung ($O(r!\cdot r^2)$,
	vgl.~\cite{epp2026subgraph}) auf logarithmischer Skala für $r=1,\ldots,12$. Bereits ab
	$r\approx 6$ übersteigt die naive Methode die rotationsbasierte Methode um mehrere
	Größenordnungen, was Lemma~\ref{lem:zustandsexplosion} und
	Bemerkung~\ref{bem:permutierbar-fazit} empirisch illustriert.}
	\label{fig:komplexitaet}
\end{figure}

\subsection{Spekulative Vorbereitung (Prefetching)}
\label{sec:prefetch}

Da der Server gemäß Definition~\ref{def:prozessautomat} für jeden Zustand $s$ die Menge
$N(s)$ der erreichbaren Folgezustände \emph{vollständig kennt}, kann er bereits während
der Bedenkzeit des Clients Vorbereitungsarbeiten für alle $s' \in N(s)$ anstoßen.

\begin{definition}[Latenzmodell]
	\label{def:latenzmodell}
	Für eine Operation $\sigma$ mit $\delta(s,\sigma)=s'$ seien:
	\begin{itemize}
		\item $T_{\mathrm{exec}}(\sigma)$: die Zeit zur Ausführung der eigentlichen
		      Fachlogik von $\sigma$ (z.\,B. Schreiboperation, Validierung),
		\item $T_{\mathrm{prep}}(s')$: die Zeit, um Ressourcen für den Zustand $s'$
		      \emph{vorzubereiten} (z.\,B. Vorab-Laden von Referenzdaten, Reservierung von
		      Inventar, Aufwärmen von Caches),
		\item $T_{\mathrm{think}}$: die Zeit zwischen dem Eintreffen der Antwort auf die
		      vorherige Operation (die in Zustand $s$ überführt hat) und dem Eintreffen
		      der nächsten Anfrage $\sigma$ des Clients (\emph{Bedenkzeit}).
	\end{itemize}
	\textbf{Ohne Prefetching} beginnt die Vorbereitung von $s'$ erst, nachdem $\sigma$
	eingetroffen ist:
	\[
		E[L_{\text{ohne}}] = T_{\mathrm{prep}}(s') + T_{\mathrm{exec}}(\sigma).
	\]
	\textbf{Mit Prefetching} startet der Server, sobald $\pi(id)=s$ erreicht ist, für
	\emph{jedes} $s'' \in N(s)$ parallel eine Vorbereitung. Trifft $\sigma$ mit
	$\delta(s,\sigma)=s'$ ein, ist die Vorbereitung für $s'$ bereits
	$\min(T_{\mathrm{think}}, T_{\mathrm{prep}}(s'))$ weit fortgeschritten:
	\[
		E[L_{\mathrm{pf}}] = \max\bigl(T_{\mathrm{prep}}(s') - T_{\mathrm{think}},\, 0\bigr)
		+ T_{\mathrm{exec}}(\sigma).
	\]
\end{definition}

\begin{satz}[Latenzgewinn durch Prefetching]
	\label{satz:latenzgewinn}
	Unter dem Latenzmodell aus Definition~\ref{def:latenzmodell} gilt für den
	Latenzgewinn $\Delta := E[L_{\text{ohne}}] - E[L_{\mathrm{pf}}]$:
	\[
		\Delta = \min\bigl(T_{\mathrm{think}},\, T_{\mathrm{prep}}(s')\bigr) \;\geq\; 0.
	\]
	Insbesondere wird die Vorbereitungszeit $T_{\mathrm{prep}}(s')$ \emph{vollständig}
	durch Prefetching verborgen, sobald $T_{\mathrm{think}} \geq T_{\mathrm{prep}}(s')$.
\end{satz}

\begin{proof}
	Wir unterscheiden zwei Fälle.

	\textbf{Fall 1: $T_{\mathrm{think}} \geq T_{\mathrm{prep}}(s')$.} Dann ist
	$T_{\mathrm{prep}}(s') - T_{\mathrm{think}} \leq 0$, also
	$\max(T_{\mathrm{prep}}(s')-T_{\mathrm{think}},0) = 0$ und somit
	$E[L_{\mathrm{pf}}] = T_{\mathrm{exec}}(\sigma)$. Es folgt
	\[
		\Delta = \bigl(T_{\mathrm{prep}}(s') + T_{\mathrm{exec}}(\sigma)\bigr)
		- T_{\mathrm{exec}}(\sigma) = T_{\mathrm{prep}}(s')
		= \min(T_{\mathrm{think}}, T_{\mathrm{prep}}(s')).
	\]

	\textbf{Fall 2: $T_{\mathrm{think}} < T_{\mathrm{prep}}(s')$.} Dann ist
	$T_{\mathrm{prep}}(s')-T_{\mathrm{think}} > 0$, also
	$\max(T_{\mathrm{prep}}(s')-T_{\mathrm{think}},0) = T_{\mathrm{prep}}(s')-T_{\mathrm{think}}$
	und somit
	\[
		E[L_{\mathrm{pf}}] = \bigl(T_{\mathrm{prep}}(s')-T_{\mathrm{think}}\bigr)
		+ T_{\mathrm{exec}}(\sigma).
	\]
	Es folgt
	\[
		\Delta = \bigl(T_{\mathrm{prep}}(s')+T_{\mathrm{exec}}(\sigma)\bigr)
		- \bigl((T_{\mathrm{prep}}(s')-T_{\mathrm{think}})+T_{\mathrm{exec}}(\sigma)\bigr)
		= T_{\mathrm{think}}
		= \min(T_{\mathrm{think}}, T_{\mathrm{prep}}(s')).
	\]

	In beiden Fällen gilt $\Delta = \min(T_{\mathrm{think}}, T_{\mathrm{prep}}(s'))$, und
	da $T_{\mathrm{think}}, T_{\mathrm{prep}}(s') \geq 0$ ist auch $\Delta \geq 0$.
	\qedhere
\end{proof}

\begin{bemerkung}[Kosten des Prefetchings]
	\label{bem:prefetch-kosten}
	Satz~\ref{satz:latenzgewinn} betrachtet nur die \emph{Latenz aus Sicht des Clients}.
	Die \emph{Kosten} des Prefetchings für den Server skalieren mit dem
	Verzweigungsgrad $d(s)=|N(s)|$ (Definition~\ref{def:prozessautomat}): Für jeden der
	$d(s)$ möglichen Folgezustände muss eine Vorbereitung angestoßen werden, von denen
	höchstens eine tatsächlich genutzt wird (\emph{spekulative Verschwendung}). Ein
	SC-REST-Server sollte daher
	\begin{itemize}
		\item den Verzweigungsgrad $d(s)$ pro Zustand durch Prozessdesign begrenzen
		      (z.\,B. $d(s) \leq D_{\max}$ für ein konfigurierbares $D_{\max}$),
		\item Vorbereitungsoperationen mit geringen Kosten $T_{\mathrm{prep}}(s')$
		      priorisieren, da nach Satz~\ref{satz:latenzgewinn} der Gewinn durch
		      $\min(T_{\mathrm{think}},T_{\mathrm{prep}}(s'))$ für kleine
		      $T_{\mathrm{prep}}(s')$ ohnehin gering ist, während die
		      \emph{verschwendeten} Ressourcen proportional zu $T_{\mathrm{prep}}(s')$
		      anfallen,
		\item idempotente bzw. abbrechbare Vorbereitungsoperationen verwenden, sodass im
		      Fall einer nicht genutzten Vorbereitung kein inkonsistenter Zustand
		      entsteht.
	\end{itemize}
	Die empirische Auswertung in Abschnitt~\ref{sec:empirisch} (Abbildung
	\ref{fig:lernkurve}, rechts) zeigt den kumulierten Effekt des Latenzgewinns über
	mehrere Prozessinstanzen unter realistischer Trefferquote.
\end{bemerkung}

\begin{figure}[h]
	\centering
	\includegraphics[width=\textwidth]{plots/plot2_latenz.pdf}
	\caption{Links: Erwartete Antwortlatenz $E[L]$ in Abhängigkeit von der Bedenkzeit
	$T_{\mathrm{think}}$ des Clients, für $T_{\mathrm{prep}}=100\,\mathrm{ms}$ und
	$T_{\mathrm{exec}}=50\,\mathrm{ms}$. Ohne Prefetching (rot, gestrichelt) ist die
	Latenz konstant $T_{\mathrm{prep}}+T_{\mathrm{exec}}=150\,\mathrm{ms}$, unabhängig von
	$T_{\mathrm{think}}$. Mit Prefetching (blau) sinkt die Latenz linear mit
	$T_{\mathrm{think}}$ bis zur unteren Schranke $T_{\mathrm{exec}}$, die bei
	$T_{\mathrm{think}}\geq T_{\mathrm{prep}}$ erreicht wird. Rechts: Latenzgewinn
	$\Delta=\min(T_{\mathrm{think}},T_{\mathrm{prep}})$ gemäß
	Satz~\ref{satz:latenzgewinn} -- ein stückweise linearer, konkaver Verlauf mit
	Sättigung bei $T_{\mathrm{prep}}$.}
	\label{fig:latenz}
\end{figure}

\subsection{Robustheit: Idempotenz und Replay-Sicherheit}
\label{sec:robustheit}

\begin{definition}[Erweiterter Automat mit Idempotenz-Schleifen]
	\label{def:idempotenz}
	Sei $P=(S,\Sigma,\delta,s_0,F)$ ein Prozessautomat. Der \emph{idempotent erweiterte
	Automat} $P^{+}=(S,\Sigma,\delta^{+},s_0,F)$ ist definiert durch
	\[
		\delta^{+}(s,\sigma) =
		\begin{cases}
			\delta(s,\sigma), & \text{falls } \delta(s,\sigma) \text{ definiert ist}, \\
			s, & \text{falls } \exists\, s_{\mathrm{prev}}\in S:\ \delta(s_{\mathrm{prev}},\sigma)=s
			      \text{ und } \sigma \text{ die zuletzt erfolgreich}\\
			   & \text{ausgeführte Operation war}, \\
			\text{undef}, & \text{sonst.}
		\end{cases}
	\]
	Das heißt: Wiederholt ein Client genau die Operation, die ihn zuletzt erfolgreich in
	den aktuellen Zustand $s$ überführt hat (z.\,B. weil die ursprüngliche Antwort durch
	einen Netzwerkfehler verloren ging), so verbleibt der Automat in $s$ und liefert
	erneut die (gecachte) Erfolgsantwort, statt eine Sequenzverletzung zu melden.
\end{definition}

\begin{satz}[Replay-Sicherheit]
	\label{satz:replay}
	Sei $P^{+}$ der idempotent erweiterte Automat zu $P$ gemäß
	Definition~\ref{def:idempotenz}, und sei $c=(c_1,\ldots,c_k)$ vollständig konform
	bezüglich $P$ mit Lauf $s_0,\ldots,s_k$. Sei $\hat c$ eine Sequenz, die aus $c$ durch
	beliebig häufige unmittelbare Wiederholung einzelner Elemente entsteht (d.\,h. jedes
	$\hat c_j \in \{c_i, c_i\}$ für ein passendes $i$, wobei die relative Reihenfolge der
	$c_i$ erhalten bleibt und Wiederholungen stets unmittelbar auf das Original folgen).
	Dann ist $\hat c$ vollständig konform bezüglich $P^{+}$, und der erreichte Endzustand
	ist identisch zu $s_k$.
\end{satz}

\begin{proof}
	Wir zeigen die Aussage durch Induktion über die Position $j$ in $\hat c$, mit der
	Invariante, dass nach Verarbeitung des durch $\hat c_1,\ldots,\hat c_j$ definierten
	Präfixes der Zustand von $P^{+}$ gleich $s_i$ ist, wobei $i$ die Anzahl der
	\emph{nicht wiederholten} (d.\,h. „neuen“) Elemente in $\hat c_1,\ldots,\hat c_j$ ist.

	\textbf{Fall „neues Element“:} Ist $\hat c_j = c_{i+1}$ ein bisher nicht in dieser
	Position wiederholtes Element, so gilt nach Induktionsannahme
	$\pi(id)=s_i$ und nach Konformität von $c$ ist $\delta(s_i,c_{i+1})=s_{i+1}$
	definiert; nach Definition~\ref{def:idempotenz} gilt $\delta^{+}=\delta$ in diesem
	Fall, also $\pi(id):=s_{i+1}$.

	\textbf{Fall „Wiederholung“:} Ist $\hat c_j$ eine unmittelbare Wiederholung von
	$\hat c_{j-1}=c_{i+1}$ (der zuletzt erfolgreich ausgeführten Operation, die
	$\pi(id)$ von $s_i$ nach $s_{i+1}$ überführt hat), so greift der zweite Fall von
	Definition~\ref{def:idempotenz} mit $s=s_{i+1}$, $s_{\mathrm{prev}}=s_i$,
	$\sigma=c_{i+1}$: $\delta^{+}(s_{i+1},c_{i+1})=s_{i+1}$, also bleibt
	$\pi(id)=s_{i+1}$ unverändert -- die Anzahl der „neuen“ Elemente bleibt $i+1$, die
	Invariante bleibt erhalten.

	Nach Verarbeitung von $\hat c$ entspricht die Anzahl der „neuen“ Elemente genau $k$
	(da $c$ vollständig $k$ Elemente hat und jede Wiederholung per Definition keinem
	neuen Element entspricht), also ist $\pi(id)=s_k$. Da $s_k\in F$ (vollständige
	Konformität von $c$), ist auch $\hat c$ vollständig konform bezüglich $P^{+}$ mit
	demselben Endzustand. \qedhere
\end{proof}

\begin{bemerkung}
	Satz~\ref{satz:replay} zeigt, dass die strikte Sequenzkontrolle aus
	Definition~\ref{def:sc-rest-server} \emph{nicht} im Widerspruch zu robusten
	Netzwerkprotokollen steht: Ein Client, der wegen eines Zeitüberschreitung (Timeout)
	eine erfolgreiche Anfrage erneut sendet, wird durch $P^{+}$ nicht fälschlich mit
	\texttt{409 Conflict} abgewiesen. In der Praxis wird $\delta^{+}$ üblicherweise mittels
	eines \emph{Idempotency-Key}-Headers realisiert (vgl. die Diskussion bei
	Helland~\cite{helland2007} zu „at-least-once“-Semantik in verteilten Systemen): Der
	Server speichert zu jedem $(id, \text{Idempotency-Key})$-Paar die zuletzt gesendete
	Antwort und liefert diese bei Wiederholung erneut aus, ohne die Fachlogik erneut
	auszuführen.
\end{bemerkung}

\section{Empirische Veranschaulichung}
\label{sec:empirisch}

Zur Veranschaulichung der formalen Ergebnisse aus Abschnitt~\ref{sec:analyse} wird die
Einführung eines SC-REST-Servers in einem simulierten Szenario betrachtet: Ein
Integrations-Team bindet schrittweise einen neuen Client an den Bestellprozess aus
Beispiel~\ref{bsp:bestellprozess} (Abbildung~\ref{fig:automat}) an.

\begin{figure}[h]
	\centering
	\includegraphics[width=\textwidth]{plots/plot4_lernkurve.pdf}
	\caption{Links: Simulierter Anteil von \texttt{409 Conflict}-Antworten über
	$n=1,\ldots,30$ bearbeitete Prozessinstanzen, modelliert als exponentieller
	Lernprozess $p(n)=p_\infty+(p_0-p_\infty)e^{-\lambda n}$ mit $p_0=0{,}42$,
	$p_\infty=0{,}015$, $\lambda=0{,}18$ (durchgezogene blaue Kurve) und einer
	binomialverteilten Stichprobe mit $m=40$ Anfragen je Instanz (rote Balken,
	Zufallsgenerator-Seed $42$). Die anfänglich hohe Fehlerrate sinkt rasch, da
	Sequenzverletzungen gemäß Satz~\ref{satz:korrektheit}(c) sofort und mit präziser
	\texttt{allowedOperations}-Angabe zurückgemeldet werden -- der Client kann sein
	Verhalten unmittelbar korrigieren. Eine Restfehlerrate $p_\infty>0$ bleibt bestehen
	(z.\,B. durch seltene Edge-Cases oder konkurrierende Prozessinstanzen). Rechts:
	Kumulierte Latenzeinsparung gemäß $\sum_n (1-p(n))\cdot m \cdot T_{\mathrm{prep}}$ mit
	$T_{\mathrm{prep}}=100\,\mathrm{ms}$ -- jede konform abgeschlossene Anfrage profitiert
	vom in Satz~\ref{satz:latenzgewinn} bewiesenen Latenzgewinn durch Prefetching, sodass
	sich der Gesamtgewinn über viele Prozessinstanzen linear mit der Trefferquote
	$1-p(n)$ aufsummiert.}
	\label{fig:lernkurve}
\end{figure}

Die Kombination der Abbildungen~\ref{fig:automat}, \ref{fig:latenz},
\ref{fig:komplexitaet} und \ref{fig:lernkurve} veranschaulicht den vollständigen
Wirkungszusammenhang des SC-REST-Patterns: Der Prozessautomat
(Abbildung~\ref{fig:automat}) definiert sowohl die zulässigen Sequenzen als auch die
Menge $\Sigma(s)$, die im Fehlerfall kommuniziert wird (Satz~\ref{satz:korrektheit});
diese Information ermöglicht serverseitiges Prefetching mit beweisbarem Latenzgewinn
(Abbildung~\ref{fig:latenz}, Satz~\ref{satz:latenzgewinn}); für lineare Abschnitte des
Prozesses bleibt die Prüfung dabei $\Theta(k)$-optimal, während permutierbare Abschnitte
strukturell auf den optimalen Subgraph Algorithmus zurückgeführt werden
(Abbildung~\ref{fig:komplexitaet}, Satz~\ref{satz:permutierbar-subgraph}); und der
Gesamtnutzen kumuliert über den Lebenszyklus einer Client-Integration
(Abbildung~\ref{fig:lernkurve}).

\section{Implementierung}
\label{sec:implementierung}

Abschnitt~\ref{sec:pattern} hat bereits die Kernfunktion \texttt{handle\_request}
(Algorithmus~\ref{lst:transition}) eingeführt. Listing~\ref{lst:flask} zeigt eine
vollständigere Referenzimplementierung als Flask-Middleware, die die
Übergangsfunktion $\delta$ deklarativ definiert, Prefetching gemäß
Abschnitt~\ref{sec:prefetch} anstößt und die Problem-Detail-Antworten aus
Abschnitt~\ref{sec:pattern} erzeugt.

\begin{lstlisting}[language=Python, label=lst:flask, caption=Referenzimplementierung eines SC-REST-Prozessautomaten in Flask]
from dataclasses import dataclass, field
from typing import Dict, Tuple, Set, List
from concurrent.futures import ThreadPoolExecutor
from flask import Flask, jsonify, request, abort

app = Flask(__name__)
executor = ThreadPoolExecutor(max_workers=8)

@dataclass
class ProcessAutomaton:
    states: Set[str]
    delta: Dict[Tuple[str, str], str]
    start: str
    final: Set[str]
    pi: Dict[str, str] = field(default_factory=dict)

    def allowed_ops(self, state: str) -> List[str]:
        return [op for (s, op) in self.delta if s == state]

    def reachable_states(self, state: str) -> Set[str]:
        return {self.delta[(state, op)]
                for op in self.allowed_ops(state)}


# Definition 2: Prozessautomat fuer den Bestellprozess
order_process = ProcessAutomaton(
    states={"s0", "s1", "s2", "s3", "s4", "s5"},
    delta={
        ("s0", "POST /cart"):        "s1",
        ("s1", "POST /cart/items"):  "s2",
        ("s2", "POST /cart/items"):  "s2",
        ("s2", "POST /checkout"):    "s3",
        ("s3", "POST /payment"):     "s4",
        ("s4", "POST /confirm"):     "s5",
    },
    start="s0",
    final={"s5"},
)


def prefetch_resources(state: str) -> None:
    """Definition 9/10: spekulative Vorbereitung fuer alle s' in N(s)."""
    for target_state in order_process.reachable_states(state):
        executor.submit(prepare_state_resources, target_state)


def prepare_state_resources(target_state: str) -> None:
    # z.B. Inventar reservieren, Zahlungsdienstleister vorab
    # autorisieren, Bestaetigungs-Template vorrendern, ...
    ...


@app.route("/orders/<order_id>/<path:op_path>", methods=["POST"])
def process_step(order_id: str, op_path: str):
    operation = f"POST /{op_path}"
    current_state = order_process.pi.get(order_id, order_process.start)

    next_state = order_process.delta.get((current_state, operation))

    if next_state is None:
        # Satz 1 (Korrektheit, Fehlerlokalisierung):
        # Zustand bleibt unveraendert, Fehler ist eindeutig lokalisiert
        return jsonify({
            "type": "https://api.example.com/probs/sequence-violation",
            "title": "Ungueltige Operation fuer aktuellen Prozesszustand",
            "status": 409,
            "instance": f"/orders/{order_id}",
            "currentState": current_state,
            "attemptedOperation": operation,
            "allowedOperations":
                order_process.allowed_ops(current_state),
        }), 409

    # Konformer Uebergang (Satz 1b: Completeness)
    result = execute_business_logic(operation, order_id)
    order_process.pi[order_id] = next_state

    # Satz 5 (Latenzgewinn): Vorbereitung fuer N(next_state)
    # waehrend der Bedenkzeit des Clients anstossen
    prefetch_resources(next_state)

    return jsonify({
        "orderId": order_id,
        "state": next_state,
        "result": result,
        "_links": hateoas_links(order_process, next_state, order_id),
    }), 201
\end{lstlisting}

\section{Von SC-REST zu HST: Horizon State Transfer}
\label{sec:sst}

\subsection{Motivation: Das Navigationsdefizit von SC-REST}
\label{sec:hst-motivation}

Das Akronym REST steht für \emph{Representational State Transfer}; namensgebend ist die
Beschränkung der \emph{Zustandslosigkeit} der Kommunikation~\cite{fielding2000}. Präzise
gefasst verlangt diese Beschränkung, dass jede Anfrage alle zur Bearbeitung notwendigen
Informationen enthält und der Server keinen \emph{Sitzungszustand} (\emph{session state})
führt -- also keinen Zustand, der ausschließlich der Kommunikation mit einem bestimmten
Client dient und ohne den frühere Anfragen dieses Clients nicht interpretierbar wären.
Davon zu unterscheiden ist der \emph{Ressourcenzustand} (in~\cite{fielding2000} auch als
Teil des \emph{application state} verstanden): persistente, über die Ressourcen-URI
adressierbare Daten, die unabhängig vom anfragenden Client existieren.

Die Bemerkung zur Zustandslosigkeit in Abschnitt~\ref{sec:pattern} stellt fest, dass
SC-REST diese Unterscheidung respektiert: $\pi(id)=s$ ist Ressourcenzustand, kein
Sitzungszustand. SC-REST ist insofern bereits REST-konform. Konsequent zu Ende gedacht
zeigt sich jedoch eine zweite, bislang unadressierte Asymmetrie zwischen Server und
Client -- das \emph{Navigationsdefizit}.

Formal kennt der Server zu jedem Zeitpunkt nicht nur den aktuellen Zustand
$\pi(id)=s$, sondern den \emph{gesamten} Prozessautomaten
$P=(S,\Sigma,\delta,s_0,F)$ -- insbesondere für jeden erreichbaren Zustand $s'\in N(s)$
wiederum $\Sigma(s')$, $N(s')$ usw. Gemäß Definition~\ref{def:sc-rest-server} übermittelt
ein SC-REST-Server in seiner Antwort jedoch nur $\Sigma(\delta(s,\sigma))$ -- also
Informationen über genau \emph{einen} Schritt voraus. Ein Client, der $k>1$ Schritte
planen möchte (etwa um zu entscheiden, ob ein mehrstufiger Vorgang überhaupt zulässig
ist, oder um -- wie in Abschnitt~\ref{sec:prefetch} diskutiert -- Ressourcen für mehrere
Schritte vorzubereiten), muss dafür $k$ aufeinanderfolgende Anfragen stellen,
\emph{obwohl} die hierfür nötige Information (die Einschränkung von $\delta$ auf den
relevanten Ausschnitt) bereits zum Zeitpunkt der ersten Antwort vollständig und
deterministisch beim Server vorliegt.

\begin{bemerkung}
	\label{bem:navigationsdefizit}
	Das Navigationsdefizit ist kein Verstoß gegen die Zustandslosigkeit im Sinne
	von~\cite{fielding2000} -- die einzelne Anfrage bleibt selbstbeschreibend, und der
	Server speichert keinen client-spezifischen Sitzungszustand. Es handelt sich jedoch
	um eine \emph{Unvollständigkeit der übertragenen Repräsentation} gegenüber dem, was
	REST mit der Beschränkung der \emph{Selbstbeschreibungsfähigkeit} und insbesondere
	mit HATEOAS anstrebt: Eine Repräsentation soll dem Client ermöglichen, ohne über die
	Anwendung hinausgehendes Wissen fortzufahren. SC-REST realisiert HATEOAS nur mit
	\emph{Horizont $h=1$} -- der folgende Abschnitt führt den Horizont $h$ als expliziten
	Parameter ein und verallgemeinert SC-REST entsprechend.
\end{bemerkung}

Ziel dieses Abschnitts ist die formale Definition eines Patterns, das dieses Defizit
behebt, ohne die Zustandslosigkeit (Abschnitt~\ref{sec:hst-zustandslosigkeit})
einzuschränken. Wir nennen dieses Pattern \emph{Horizon State Transfer (HST)}, da --
wie sich zeigen wird -- die zusätzlich übertragene Information genau der von $s$ aus in
höchstens $h$ Schritten erreichbare Teilgraph des Prozessautomaten ist
(formal behandelt in~\cite{epp2026subgraph}), also der \emph{Navigationshorizont} der
Tiefe $h$. Der Name ist bewusst in Analogie zu \emph{Representational State Transfer}
gewählt: Während REST in jeder Antwort eine \emph{Repräsentation} des aktuellen Zustands
überträgt, überträgt HST zusätzlich den vollständigen \emph{Horizont} $G_s^h$ des
Zustandsraums -- den lokal relevanten Ausschnitt des Navigationswissens, das der Server
ohnehin vollständig besitzt, begrenzt durch den Parameter $h$.

\subsection{Das HST-Pattern: Formale Definition}
\label{sec:hst-definition}

\begin{definition}[$h$-Schritt-Horizont-Subgraph]
	\label{def:horizont-subgraph}
	Sei $P=(S,\Sigma,\delta,s_0,F)$ ein Prozessautomat, $s\in S$ und $h\in\mathbb{N}_0$.
	Wir fassen $P$ als gerichteten, mit $\Sigma$ beschrifteten Graphen $(S,\to)$ auf, mit
	$s_1\xrightarrow{\sigma}s_2 :\Leftrightarrow \delta(s_1,\sigma)=s_2$. Für
	$(c_1,\ldots,c_k)\in\Sigma^k$ sei $\delta^{*}(s,(c_1,\ldots,c_k))$ der gemäß
	Definition~\ref{def:lauf} erreichte Zustand nach Ausführung von
	$(c_1,\ldots,c_k)$ ab $s$, mit $\delta^{*}(s,())=s$. Der \emph{$h$-Schritt-Horizont-Subgraph}
	$G_s^h=(V_s^h,E_s^h)$ ist definiert durch
	\[
		V_s^h := \bigl\{\, s' \in S \;\bigm|\; \exists\, k\le h,\ \exists\, c\in\Sigma^k:\
		s'=\delta^{*}(s,c) \,\bigr\}
	\]
	und, für $h\ge1$,
	\[
		E_s^h := \bigl\{\, (s_1,\sigma,s_2) \;\bigm|\; s_1\in V_s^{h-1},\ \sigma\in\Sigma(s_1),\
		s_2=\delta(s_1,\sigma) \,\bigr\},
	\]
	mit $E_s^0:=\emptyset$. Jeder Knoten $s'\in V_s^h\setminus V_s^{h-1}$ (für $h\ge1$)
	heißt \emph{Randknoten} von $G_s^h$; die Menge der Randknoten wird mit $\partial
	G_s^h$ bezeichnet.
\end{definition}

Anschaulich ist $G_s^h$ der induzierte Teilgraph, der durch eine Breitensuche (BFS) der
Tiefe $h$ ab $s$ im Prozessautomaten entsteht -- einschließlich aller Kanten, deren
Startknoten in höchstens $h-1$ Schritten ab $s$ erreichbar ist.

\begin{bemerkung}
	\label{bem:sc-rest-spezialfall}
	$G_s^0=(\{s\},\emptyset)$, und $G_s^1$ entspricht genau der in
	Definition~\ref{def:sc-rest-server} übertragenen Information: $V_s^1=\{s\}\cup N(s)$
	mit den Kanten $\{(s,\sigma,\delta(s,\sigma))\mid \sigma\in\Sigma(s)\}$. \textbf{SC-REST
	ist somit der Spezialfall $h=1$ von HST.}
\end{bemerkung}

\begin{definition}[HST-Repräsentation]
	\label{def:hst-repraesentation}
	Sei $(id,s)$ eine Prozessinstanz gemäß Definition~\ref{def:prozessinstanz} und
	$h\in\mathbb{N}$ ein (server- oder anwendungsseitig konfigurierter)
	\emph{Horizontparameter}. Die \emph{HST-Repräsentation} der Prozessinstanz ist das
	Tupel
	\[
		R(id,s,h) := \bigl(\, \rho(id,s),\ s,\ G_s^h \,\bigr),
	\]
	wobei $\rho(id,s)$ die fachliche Ressourcenrepräsentation (Nutzdaten der
	Prozessinstanz) bezeichnet und $G_s^h$ gemäß Definition~\ref{def:horizont-subgraph}
	als annotierter Graph übertragen wird: Jeder Knoten $s'\in V_s^h$ wird auf eine URI
	(bzw.\ einen URI-Template-Parameter) abgebildet, jede Kante $(s_1,\sigma,s_2)\in
	E_s^h$ auf einen HATEOAS-Link mit Methode, Pfad und annotiertem Zielzustand $s_2$.
\end{definition}

\begin{definition}[HST-Operationssemantik]
	\label{def:hst-server}
	Sei $P=(S,\Sigma,\delta,s_0,F)$ ein Prozessautomat, $\pi(id)=s$ der aktuelle Zustand
	einer Prozessinstanz und $h\in\mathbb{N}$ der Horizontparameter. Auf eine eingehende
	Anfrage $\sigma\in\Sigma$ reagiert ein \emph{HST-Server} wie folgt:
	\begin{enumerate}
		\item \textbf{Falls $\delta(s,\sigma)$ definiert ist:} Wie in
		Definition~\ref{def:sc-rest-server} (Fachlogik ausführen, $\pi(id):=
		\delta(s,\sigma)$), jedoch antwortet der Server mit der HST-Repräsentation
		$R(id,\delta(s,\sigma),h)$ statt nur mit $\Sigma(\delta(s,\sigma))$.
		\item \textbf{Falls $\delta(s,\sigma)$ undefiniert ist:} Wie in
		Definition~\ref{def:sc-rest-server} (\texttt{409 Conflict} mit
		Problem-Detail-Objekt, $\pi(id)$ bleibt unverändert); zusätzlich enthält die
		Antwort weiterhin $G_s^h$ (unverändert, da $\pi(id)$ unverändert bleibt), sodass
		der Client den Sequenzfehler unmittelbar im Kontext seines bereits bekannten
		Horizont-Subgraphen einordnen kann.
	\end{enumerate}
\end{definition}

\begin{beispiel}[Bestellprozess mit Horizont $h=2$]
	\label{bsp:hst-bestellprozess}
	Für den Prozessautomaten aus Beispiel~\ref{bsp:bestellprozess} und $s=s_1$ (Zustand
	nach \texttt{POST /cart}) ergibt sich für $h=2$:
	\[
		V_{s_1}^2 = \{s_1,s_2,s_3\},
	\]
	\[
		\begin{aligned}
			E_{s_1}^2 = \bigl\{\, &(s_1,\texttt{POST /cart/items},s_2),\
				(s_2,\texttt{POST /cart/items},s_2),\\
				&(s_2,\texttt{POST /checkout},s_3) \,\bigr\}.
		\end{aligned}
	\]
	Ein Client, der $G_{s_1}^2$ erhält, kann \emph{ohne weitere Serveranfrage}
	feststellen, dass \texttt{POST /checkout} aus $s_1$ direkt nicht zulässig ist, jedoch
	nach genau einem \texttt{POST /cart/items} zulässig wird -- eine Information, die in
	der SC-REST-Repräsentation ($h=1$, nur $\Sigma(s_1)=\{\texttt{POST /cart/items}\}$)
	nicht enthalten ist.
\end{beispiel}

\subsection{Zustandslosigkeit von HST}
\label{sec:hst-zustandslosigkeit}

\begin{satz}[Zustandslosigkeit von HST]
	\label{satz:hst-zustandslosigkeit}
	Für jeden Horizont $h\in\mathbb{N}$ erfüllt ein HST-Server gemäß
	Definition~\ref{def:hst-server} die Zustandslosigkeits-Beschränkung
	nach~\cite{fielding2000}: Der Server führt keinen client-spezifischen
	Sitzungszustand, und jede Antwort ist eine deterministische Funktion ausschließlich
	von (a) dem Prozessautomaten $P$ und (b) dem Ressourcenzustand $\pi(id)=s$.
\end{satz}

\begin{proof}
	Wir zeigen, dass $R(id,s,h)$ gemäß Definition~\ref{def:hst-repraesentation}
	vollständig durch $(P,id,s,h)$ bestimmt ist, unabhängig davon, welcher Client die
	Anfrage stellt oder welche Anfragen zuvor gestellt wurden.

	(1) $\rho(id,s)$ ist nach Definition~\ref{def:prozessinstanz} die fachliche
	Repräsentation der Ressource $id$ im Zustand $s$ -- eine Funktion von $(id,s)$,
	unabhängig vom anfragenden Client.

	(2) $G_s^h$ ist nach Definition~\ref{def:horizont-subgraph} eine Funktion von
	$(P,s,h)$: Die Mengen $V_s^h$ und $E_s^h$ werden ausschließlich mittels
	$\delta^{*}$ -- also iterierter Anwendung von $\delta$, einer Komponente von $P$ --
	aus $s$ berechnet. In diese Berechnung geht insbesondere keine Information über die
	Identität des anfragenden Clients, über dessen vorherige Anfragen oder über andere
	Prozessinstanzen ein.

	(3) $h$ ist nach Voraussetzung ein server- oder anwendungsseitig (nicht
	client-sitzungsseitig) konfigurierter Parameter -- etwa eine Konstante des
	Prozessautomaten $P$ oder ein in der aktuellen Anfrage selbst übermittelter Parameter
	(z.\,B.\ Query-Parameter \texttt{?horizon=2}); in beiden Fällen ist $h$ Teil der
	Eingabe der aktuellen Anfrage und keine gespeicherte Information.

	Damit ist $R(id,s,h)=(\rho(id,s),s,G_s^h)$ eine Funktion von $(P,id,s,h)$, wobei $P$
	statisch (bzw.\ unabhängig versioniert, siehe Abschnitt~\ref{sec:hst-cache}) und
	$\pi(id)=s$ Ressourcenzustand im Sinne der Bemerkung zur Zustandslosigkeit in
	Abschnitt~\ref{sec:pattern} ist. Insbesondere benötigt der Server zur Beantwortung
	\emph{keinen} zusätzlichen Speicher pro Client oder pro Sitzung -- die einzige
	persistente, über die Anfrage hinaus gespeicherte Information ist $\pi(id)$, welche
	bereits für SC-REST ($h=1$) als Ressourcenzustand klassifiziert wurde. Da diese
	Klassifikation nicht von $h$ abhängt, gilt sie für beliebiges $h\in\mathbb{N}$, also
	auch für HST. \qedhere
\end{proof}

\begin{bemerkung}
	\label{bem:zustandslosigkeit-frei}
	Satz~\ref{satz:hst-zustandslosigkeit} zeigt, dass die Erhöhung des Horizonts $h$
	\emph{zustandslosigkeitsneutral} ist: $h$ beeinflusst ausschließlich die
	\emph{Größe} der Antwort (Abschnitt~\ref{sec:hst-groesse}), nicht deren Charakter als
	zustandslose, selbstbeschreibende Nachricht. Der vermeintliche Zielkonflikt zwischen
	„mehr Navigationsinformation für den Client“ und „REST-Konformität“ -- der die in
	Bemerkung~\ref{bem:navigationsdefizit} beschriebene Beschränkung auf $h=1$ in der
	Praxis oft motiviert -- existiert formal nicht; lediglich der in
	Abschnitt~\ref{sec:hst-kosten} formalisierte Kosten-Nutzen-Tradeoff begrenzt $h$ auf
	einen endlichen, optimalen Wert $h^\ast$.
\end{bemerkung}
\subsection{Größe des Horizont-Subgraphen}
\label{sec:hst-groesse}

\begin{lemma}[Größenschranke für $G_s^h$]
	\label{lem:horizont-groesse}
	Sei $P=(S,\Sigma,\delta,s_0,F)$ ein Prozessautomat und $d:=\max_{s'\in S}d(s')$ der
	maximale Verzweigungsgrad (Definition~\ref{def:prozessautomat}). Für jeden Zustand
	$s\in S$ und jeden Horizont $h\in\mathbb{N}_0$ gilt
	\[
		|V_s^h| \;\le\; \sum_{i=0}^{h} d^{\,i} \;=\;
		\begin{cases}
			h+1, & d=1,\\[2pt]
			\dfrac{d^{h+1}-1}{d-1}, & d\ne1.
		\end{cases}
	\]
	Gleichheit gilt genau dann, wenn der von $s$ aus erreichbare Teil von $P$ bis zur
	Tiefe $h$ ein Baum mit konstantem Verzweigungsgrad $d$ ist, d.\,h.\ wenn
	$\delta^{*}$ auf Pfaden der Länge $\le h$ ab $s$ injektiv ist (keine zwei
	verschiedenen solchen Pfade führen zum selben Zustand) und jeder Zustand auf Tiefe
	$<h$ genau $d$ ausgehende Operationen besitzt.
\end{lemma}

\begin{proof}
	Wir zerlegen $V_s^h$ in Schichten $L_0,\ldots,L_h$ mit $L_0:=\{s\}$ und, für
	$i\ge1$, $L_i:=V_s^i\setminus V_s^{i-1}$ (mit $V_s^{-1}:=\emptyset$). Wir zeigen
	$|L_i|\le d^i$ für alle $i\le h$ durch Induktion über $i$.

	\emph{Induktionsanfang} ($i=0$): $L_0=\{s\}$, also $|L_0|=1=d^0$.

	\emph{Induktionsschritt}: Sei $|L_{i-1}|\le d^{i-1}$ gezeigt. Jeder Knoten
	$s'\in L_i$ ist nach Definition~\ref{def:horizont-subgraph} von der Form
	$s'=\delta(s'',\sigma)$ für ein $s''\in L_{i-1}$ und $\sigma\in\Sigma(s'')$ -- denn
	$s'\in V_s^i\setminus V_s^{i-1}$ bedeutet, dass der kürzeste Pfad von $s$ nach $s'$
	genau Länge $i$ hat, also über einen Knoten $s''$ auf Tiefe $i-1$ verläuft
	(andernfalls wäre $s'$ bereits über einen kürzeren Pfad in $V_s^{i-1}$ enthalten).
	Da $d(s'')\le d$ für alle $s''\in S$ (Definition von $d$ als Maximum), gibt es zu
	jedem $s''\in L_{i-1}$ höchstens $d$ Operationen $\sigma\in\Sigma(s'')$ und damit
	höchstens $d$ Kandidaten für Nachfolgeknoten. Da
	$L_i\subseteq\bigcup_{s''\in L_{i-1}}\{\delta(s'',\sigma)\mid\sigma\in\Sigma(s'')\}$,
	folgt
	\[
		|L_i| \le |L_{i-1}|\cdot d \le d^{i-1}\cdot d = d^i.
	\]

	Damit folgt $|V_s^h|=\sum_{i=0}^{h}|L_i|\le\sum_{i=0}^{h}d^i$.

	\emph{Gleichheitsfall}: Gleichheit $|L_i|=d\cdot|L_{i-1}|$ für alle $i\le h$ gilt
	genau dann, wenn (a) jeder Knoten $s''\in L_{i-1}$ \emph{genau} $d$ ausgehende
	Operationen besitzt (der maximale Verzweigungsgrad wird an jedem dieser Knoten
	tatsächlich erreicht) \emph{und} (b) die resultierenden $d\cdot|L_{i-1}|$
	Nachfolgeknoten paarweise verschieden sind und nicht bereits in $V_s^{i-1}$ liegen.
	Bedingungen (a) und (b) für alle $i\le h$ sind äquivalent dazu, dass $\delta^{*}$ auf
	Pfaden der Länge $\le h$ ab $s$ injektiv ist und jeder Knoten auf Tiefe $<h$
	Verzweigungsgrad genau $d$ hat -- d.\,h.\ der von $s$ aus erreichbare Teilgraph bis
	Tiefe $h$ ist ein $d$-närer Baum. \qedhere
\end{proof}

\begin{bemerkung}
	\label{bem:horizont-groesse-praxis}
	Lemma~\ref{lem:horizont-groesse} zeigt, dass $|V_s^h|$ für $d\ge2$
	\emph{exponentiell} in $h$ wächst. Für die in Abschnitt~\ref{sec:permutierbar}
	betrachteten Sammelzustände mit $r$ permutierbaren optionalen Teilschritten ist dies
	kein Widerspruch zu Bemerkung~\ref{bem:permutierbar-fazit}: Dort wird $|S|\ge2^r$ als
	\emph{untere} Schranke für die naive Kodierung des \emph{gesamten} Automaten gezeigt,
	während Lemma~\ref{lem:horizont-groesse} eine \emph{obere} Schranke für den
	\emph{lokal pro Anfrage übertragenen} Teilgraphen $G_s^h$ liefert -- für $h<r$ ist
	$|V_s^h|$ um Größenordnungen kleiner als $|S|$. Abbildung~\ref{fig:horizont-groesse}
	illustriert das Wachstum von $|V_s^h|$ für $d\in\{1,2,3\}$.
\end{bemerkung}

\begin{figure}[ht]
	\centering
	\begin{tikzpicture}[
		every node/.style={circle, draw, minimum size=7mm, font=\footnotesize},
		level distance=13mm,
		level 1/.style={sibling distance=30mm},
		level 2/.style={sibling distance=15mm},
		edge from parent/.style={draw, -{Stealth[length=2mm]}},
		>=Stealth
		]
		\node[fill=mainblue!25, draw=mainblue, line width=0.9pt] (s) {$s$}
			child { node[fill=mainblue!25, draw=mainblue, line width=0.9pt] (a) {$s_a$}
				child { node[fill=mainblue!25, draw=mainblue, line width=0.9pt] {$s_{a1}$}
					child { node[draw=lightgray, text=darkgray, dashed] {$\cdots$} edge from parent[draw=lightgray, dashed] }
				}
				child { node[fill=mainblue!25, draw=mainblue, line width=0.9pt] {$s_{a2}$}
					child { node[draw=lightgray, text=darkgray, dashed] {$\cdots$} edge from parent[draw=lightgray, dashed] }
				}
			}
			child { node[fill=mainblue!25, draw=mainblue, line width=0.9pt] (b) {$s_b$}
				child { node[fill=mainblue!25, draw=mainblue, line width=0.9pt] {$s_{b1}$}
					child { node[draw=lightgray, text=darkgray, dashed] {$\cdots$} edge from parent[draw=lightgray, dashed] }
				}
				child { node[fill=mainblue!25, draw=mainblue, line width=0.9pt] {$s_{b2}$}
					child { node[draw=lightgray, text=darkgray, dashed] {$\cdots$} edge from parent[draw=lightgray, dashed] }
				}
			};
		\node[draw=none, fill=none, font=\scriptsize] at ($(s)!0.4!(a) + (-0.25,0)$) {$\sigma_1$};
		\node[draw=none, fill=none, font=\scriptsize] at ($(s)!0.4!(b) + (0.25,0)$) {$\sigma_2$};
		\draw[dashed, accentred, line width=0.8pt, rounded corners]
			($(s.north)+(-4.6,0.35)$) rectangle ($(b.south)+(4.6,-1.55)$);
		\node[draw=none, fill=none, shape=rectangle, accentred, font=\scriptsize, anchor=north west]
			at ($(s.north)+(-4.55,0.3)$) {Horizont $h=2$: $G_s^2$};
	\end{tikzpicture}
	\caption{Schematische Darstellung des Horizont-Subgraphen $G_s^2$ (blau, innerhalb
	der gestrichelten roten Umrandung) für einen Prozessautomaten mit Verzweigungsgrad
	$d=2$ ab dem aktuellen Zustand $s$. Gemäß Lemma~\ref{lem:horizont-groesse} enthält
	$G_s^2$ höchstens $\sum_{i=0}^2 2^i=1+2+4=7$ Knoten -- hier mit Gleichheit, da der
	dargestellte Ausschnitt ein binärer Baum ist. Die hellgrau gestrichelten Knoten
	($\cdots$) liegen auf Tiefe $3$, also außerhalb des Horizonts, und sind dem Client
	ohne weitere Anfrage nicht bekannt: Erst nach einem Zustandsübergang innerhalb des
	Horizonts (z.\,B.\ $s\xrightarrow{\sigma_1}s_a$) liefert der Server den
	\emph{aktualisierten} Subgraphen $G_{s_a}^2$, dessen Randknoten dann diese zuvor
	verborgenen Zustände enthält. Der Horizont „wandert“ also mit der Prozessinstanz
	mit, ohne dass der Server jemals den vollständigen Automaten $P$ in einer einzelnen
	Antwort übertragen muss.}
	\label{fig:horizont-baum}
\end{figure}

\begin{figure}[ht]
	\centering
	\includegraphics[width=\textwidth]{plots/plot5_horizont_groesse.pdf}
	\caption{Links: Anzahl der Knoten $|V_s^h|$ des Horizont-Subgraphen (logarithmische
	Skala) in Abhängigkeit vom Horizont $h$, für maximale Verzweigungsgrade $d=1$
	(lineare Prozessabschnitte ohne Verzweigung, z.\,B.\ der Hauptpfad in
	Beispiel~\ref{bsp:bestellprozess}), $d=2$ und $d=3$, gemäß der Schranke aus
	Lemma~\ref{lem:horizont-groesse} (mit angenommener Gleichheit als Worst-Case). Für
	$d=1$ wächst $|V_s^h|$ nur linear ($h+1$), für $d\ge2$ exponentiell -- bereits für
	$d=3$ und $h=8$ ergeben sich knapp $10^4$ Knoten. Rechts: geschätzte
	Übertragungsgröße der HST-Repräsentation $R(id,s,h)$ in Kilobyte, unter der
	konservativen Annahme von ca.\ 250\,Byte pro Knoten bzw.\ Kante (URI, Methode,
	Zustandskennung, JSON-Strukturoverhead). Die gestrichelte Linie markiert 8\,KB als
	groben Richtwert für eine „leichtgewichtige“ HTTP-Antwort. Für $d=2$ wird dieser Wert
	bereits bei $h\approx4$, für $d=3$ bei $h\approx3$ überschritten -- dies motiviert
	den in Abschnitt~\ref{sec:hst-kosten} formalisierten kostenoptimalen Horizont
	$h^\ast$, der für praxisrelevante Parameter deutlich kleiner ist als die in
	Abschnitt~\ref{sec:permutierbar} betrachteten Werte $r\le20$.}
	\label{fig:horizont-groesse}
\end{figure}

\subsection{Lokale Vollständigkeit und Planungshorizont}
\label{sec:hst-vollstaendigkeit}

\begin{satz}[Lokale Vollständigkeit von $G_s^h$]
	\label{satz:lokale-vollstaendigkeit}
	Sei $(id,s)$ eine Prozessinstanz und $G_s^h$ der zugehörige Horizont-Subgraph gemäß
	Definition~\ref{def:horizont-subgraph}. Für jede Aufrufsequenz $c=(c_1,\ldots,c_k)$
	mit $k\le h$ gilt: Die Frage, ob $c$ ab $s$ (partiell) konform ist
	(Definition~\ref{def:konformitaet}), und -- falls ja -- welcher Zustand
	$\delta^{*}(s,c)$ erreicht wird, ist \emph{vollständig aus $G_s^h$ entscheidbar},
	ohne weitere Anfrage an den Server.
\end{satz}

\begin{proof}
	Wir zeigen zunächst die monotone Inklusion $V_s^i\subseteq V_s^j$ für $i\le j$: Jeder
	in höchstens $i$ Schritten ab $s$ erreichbare Zustand ist insbesondere auch in
	höchstens $j\ge i$ Schritten erreichbar (man kann dieselbe Sequenz verwenden), also
	folgt die Inklusion direkt aus Definition~\ref{def:horizont-subgraph}.

	Wir zeigen die Behauptung nun durch Induktion über $k\le h$.

	\emph{Induktionsanfang} ($k=0$): $c=()$ ist nach Definition~\ref{def:konformitaet}
	trivial konform mit $\delta^{*}(s,())=s\in V_s^h$ für jedes $h\ge0$.

	\emph{Induktionsschritt}: Sei die Aussage für $k-1<h$ Schritte gezeigt, d.\,h.\ es
	existiert (unter der Annahme, dass $(c_1,\ldots,c_{k-1})$ konform ist) ein Zustand
	$s_{k-1}:=\delta^{*}(s,(c_1,\ldots,c_{k-1}))\in V_s^{k-1}$, und dieser Zustand sowie
	die Konformität von $(c_1,\ldots,c_{k-1})$ sind aus $G_s^h$ bestimmbar. Wegen
	$k-1\le h-1$ und der oben gezeigten Monotonie gilt $s_{k-1}\in V_s^{h-1}$.

	Zu prüfen ist, ob $\delta(s_{k-1},c_k)$ definiert ist. Da $s_{k-1}\in V_s^{h-1}$, ist
	nach Definition~\ref{def:horizont-subgraph} \emph{jede} Operation
	$\sigma\in\Sigma(s_{k-1})$ Quelle einer Kante
	$(s_{k-1},\sigma,\delta(s_{k-1},\sigma))\in E_s^h$ -- die Kantenmenge $E_s^h$
	enthält per Definition alle ausgehenden Kanten \emph{jedes} Knotens aus
	$V_s^{h-1}$, unabhängig davon, ob deren Zielknoten neu (Randknoten von $G_s^h$)
	oder bereits in $V_s^{h-1}$ enthalten ist. Somit gilt:
	\begin{itemize}
		\item $\delta(s_{k-1},c_k)$ ist definiert $\iff$ es existiert eine Kante
		      $(s_{k-1},c_k,s_k)\in E_s^h$ für ein $s_k\in V_s^h$ -- diese Bedingung ist
		      direkt aus der (endlichen) Kantenmenge $E_s^h$ ablesbar.
		\item Ist $\delta(s_{k-1},c_k)$ definiert, so ist der gemäß
		      Definition~\ref{def:prozessautomat} eindeutige Zielknoten $s_k$ dieser
		      Kante genau $\delta^{*}(s,(c_1,\ldots,c_k))$, und $s_k\in V_s^{k}\subseteq
		      V_s^h$ (wieder nach Monotonie, $k\le h$).
		\item Ist $\delta(s_{k-1},c_k)$ undefiniert, so ist $c$ ab Schritt $k$ nicht
		      konform -- auch diese Information ist aus $E_s^h$ ablesbar (Abwesenheit
		      einer entsprechenden Kante mit Quelle $s_{k-1}$ und Beschriftung $c_k$).
	\end{itemize}
	In beiden Fällen ist die Konformität von $c$ bis Schritt $k$ und ggf.\ der erreichte
	Zustand $\delta^{*}(s,(c_1,\ldots,c_k))$ vollständig aus $G_s^h$ bestimmbar. Per
	Induktion folgt die Behauptung für alle $k\le h$. \qedhere
\end{proof}

\begin{bemerkung}
	\label{bem:hst-permutierbar}
	Satz~\ref{satz:lokale-vollstaendigkeit} hat eine direkte Konsequenz für die in
	Abschnitt~\ref{sec:permutierbar} behandelten Sammelzustände: Ist $s_c$ ein
	Sammelzustand mit $r$ permutierbaren optionalen Teilschritten
	(Definition~\ref{def:sammelzustand}) und wählt der Server den Horizont $h\ge r$, so
	enthält $G_{s_c}^h$ -- mit $d=2$ an jedem Zwischenzustand (Operation $\sigma_j$
	ausführen oder nicht) und damit nach Lemma~\ref{lem:horizont-groesse} mit
	$|V_{s_c}^h|\le\sum_{i=0}^r 2^i=2^{r+1}-1$ Knoten -- den \emph{vollständigen}
	Entscheidungsbaum aller $r!$ möglichen Ausführungsreihenfolgen. Ein Client kann damit
	für eine \emph{beliebige}, lokal geplante Reihenfolge $c'$ der $r$ Teilschritte
	mittels Satz~\ref{satz:lokale-vollstaendigkeit} \emph{vorab und lokal} prüfen, ob
	$c'$ konform ist -- und kann zusätzlich, sofern $c'$ als Beobachtungsfolge im Sinne
	von Definition~\ref{def:signaturfolge-sammel} interpretiert wird, den Subgraph
	Algorithmus aus~\cite{epp2026subgraph} \emph{client-seitig} auf den in $G_{s_c}^h$
	enthaltenen Signaturen ausführen (Satz~\ref{satz:permutierbar-subgraph}(b),
	$\Theta(r^3)$). Die Konformitätsprüfung für Sammelzustände wird damit
	\emph{vollständig dezentralisiert}: Der Server muss sie nur noch im Fehlerfall (echte
	Sequenzverletzung, etwa durch einen vom Plan abweichenden Client) erneut
	durchführen, vgl.\ Definition~\ref{def:hst-server}(2).
\end{bemerkung}
\subsection{Kostenoptimaler Horizont}
\label{sec:hst-kosten}

\begin{definition}[Kostenmodell für den Horizont]
	\label{def:kostenmodell-horizont}
	Sei $\alpha>0$ die marginale Ãœbertragungskosten pro Knoten/Kante des
	Horizont-Subgraphen (z.\,B.\ in normierten Bandbreiten- oder Byte-Einheiten) und
	$\beta>0$ die Kosten eines zusätzlichen Anfrage-Antwort-Zyklus
	(\emph{Round-Trip}), der anfällt, wenn der Client eine Operation ausführen möchte,
	deren Zielzustand außerhalb von $G_s^h$ liegt (\emph{Horizontüberschreitung}). Wir
	modellieren die Anzahl der vom Client über den Horizont $h$ hinaus benötigten
	Schritte als eine Zufallsvariable mit geometrisch fallender Ãœberschreitungswahrscheinlichkeit:
	Mit \emph{Fortsetzungswahrscheinlichkeit} $p\in(0,1)$ sei
	$\Pr[\text{Client überschreitet Horizont }h] = p^{h+1}$, und im Überschreitungsfall
	seien im Erwartungswert $\frac{1}{1-p}$ zusätzliche Round-Trips erforderlich
	(geometrische Reihe über weitere mögliche Überschreitungen). Die \emph{erwarteten
	Gesamtkosten} für Horizont $h\in\mathbb{N}_0$ sind dann
	\[
		C(h) := \alpha\cdot|V_s^h| + \beta\cdot\frac{p^{h+1}}{1-p}
		     \;=\; \alpha\sum_{i=0}^{h} d^i \;+\; \beta\cdot\frac{p^{h+1}}{1-p},
	\]
	wobei $|V_s^h|=\sum_{i=0}^h d^i$ gemäß Lemma~\ref{lem:horizont-groesse}
	(Gleichheitsfall, als Worst-Case-Schätzung des Verzweigungsgrads $d\ge1$) verwendet
	wird.
\end{definition}

\begin{satz}[Existenz und Charakterisierung des optimalen Horizonts]
	\label{satz:optimaler-horizont}
	Sei $C(h)$ wie in Definition~\ref{def:kostenmodell-horizont}, mit $d\ge1$,
	$p\in(0,1)$, $\alpha,\beta>0$. Dann gilt:
	\begin{enumerate}
		\item[(a)] Die Differenzfunktion $\Delta C(h):=C(h+1)-C(h)=\alpha\,d^{h+1} -
		\beta\,p^{h+1}$ ist streng monoton wachsend in $h\in\mathbb{N}_0$.
		\item[(b)] $C$ ist \emph{unimodal}: Es existiert ein eindeutiges
		$h^\ast\in\mathbb{N}_0$, sodass $C(0)>C(1)>\cdots>C(h^\ast)\le
		C(h^\ast+1)<C(h^\ast+2)<\cdots$ (mit der Konvention, dass die linke Kette leer
		ist, falls bereits $\Delta C(0)\ge0$). Insbesondere gilt
		$C(h^\ast)=\min_{h\in\mathbb{N}_0}C(h)$.
		\item[(c)] Es gilt die geschlossene Form
		\[
			h^\ast = \max\!\left(0,\ \left\lceil \log_{d/p}\!\left(\frac{\beta}{\alpha}\right)
			\right\rceil - 1\right).
		\]
	\end{enumerate}
\end{satz}

\begin{proof}
	\textbf{(a)} Da $d\ge1$ und $0<p<1$, gilt $d/p>1$ (im Fall $d=1$ ist $d/p=1/p>1$
	wegen $p<1$; im Fall $d>1$ erst recht). Für die zweite Differenz gilt
	\[
		\Delta C(h+1)-\Delta C(h)
		= \bigl(\alpha d^{h+2}-\beta p^{h+2}\bigr) - \bigl(\alpha d^{h+1}-\beta p^{h+1}\bigr)
		= \alpha d^{h+1}(d-1) + \beta p^{h+1}(1-p).
	\]
	Da $d\ge1$ (also $d-1\ge0$), $0<p<1$ (also $1-p>0$) und $\alpha,\beta>0$, ist der
	zweite Summand $\beta p^{h+1}(1-p)$ stets strikt positiv und der erste Summand
	nicht-negativ. Also $\Delta C(h+1)-\Delta C(h)>0$ für alle $h\ge0$, d.\,h.\ $\Delta C$
	ist streng monoton wachsend.

	\textbf{(b)} Für $h\to\infty$ gilt $\Delta C(h)=\alpha d^{h+1}-\beta p^{h+1}$. Da
	$p^{h+1}\to0$ und $d^{h+1}\ge1$ für alle $h\ge0$ (mit $d^{h+1}\to\infty$ falls $d>1$
	bzw.\ $d^{h+1}=1$ konstant falls $d=1$), folgt $\Delta C(h)\to\alpha\cdot
	\lim_{h\to\infty}d^{h+1}\ge\alpha>0$ für $h\to\infty$ -- insbesondere wird
	$\Delta C(h)$ für hinreichend großes $h$ positiv. Da $\Delta C$ nach (a) streng
	monoton wachsend ist und für große $h$ positiv wird, gibt es höchstens einen
	Vorzeichenwechsel von „$<0$“ zu „$\ge0$“. Sei
	\[
		h^\ast := \min\{\, h\in\mathbb{N}_0 \mid \Delta C(h)\ge0 \,\}
	\]
	(wohldefiniert, da die Menge nichtleer ist nach obigem Grenzwertargument). Für
	$h<h^\ast$ gilt $\Delta C(h)<0$ (Minimalität von $h^\ast$ und Monotonie von
	$\Delta C$ schließen $\Delta C(h)\ge0$ für $h<h^\ast$ aus), also
	$C(h+1)<C(h)$: $C$ ist auf $\{0,\ldots,h^\ast\}$ streng fallend. Für $h\ge h^\ast$
	gilt $\Delta C(h)\ge0$ wegen Monotonie von $\Delta C$ und $\Delta C(h^\ast)\ge0$,
	also $C(h+1)\ge C(h)$, mit strikter Ungleichung für $h>h^\ast$ (da $\Delta C(h)>
	\Delta C(h^\ast)\ge0$ für $h>h^\ast$ nach strikter Monotonie). Also ist $C$ auf
	$\{h^\ast,h^\ast+1,\ldots\}$ streng wachsend. Damit ist $C$ unimodal mit eindeutigem
	Minimum bei $h^\ast$.

	\textbf{(c)} Die definierende Bedingung $\Delta C(h)\ge0$ ist äquivalent zu
	\[
		\alpha\,d^{h+1}\ge\beta\,p^{h+1}
		\iff \left(\frac{d}{p}\right)^{h+1}\ge\frac{\beta}{\alpha}
		\iff (h+1)\log\frac{d}{p}\ge\log\frac{\beta}{\alpha}
		\iff h+1\ge\log_{d/p}\frac{\beta}{\alpha},
	\]
	wobei die letzte Äquivalenz gilt, da $\log(d/p)>0$ nach (a) und das Logarithmieren
	einer Ungleichung mit positiver Basis $>1$ die Richtung erhält. Der kleinste
	ganzzahlige Wert $h\ge0$ mit $h+1\ge\log_{d/p}(\beta/\alpha)$ ist
	\[
		h = \max\!\left(0,\ \left\lceil\log_{d/p}\!\left(\frac{\beta}{\alpha}\right)\right\rceil-1\right)
	\]
	(die Maximumsbildung mit $0$ deckt den Fall ab, dass bereits $\Delta C(0)\ge0$ gilt,
	also $\log_{d/p}(\beta/\alpha)\le1$). Nach Definition von $h^\ast$ in (b) ist dies
	genau $h^\ast$. \qedhere
\end{proof}

\begin{bemerkung}
	\label{bem:optimaler-horizont-interpretation}
	Satz~\ref{satz:optimaler-horizont} bestätigt formal die in
	Bemerkung~\ref{bem:zustandslosigkeit-frei} aufgeworfene Beobachtung: Der optimale
	Horizont $h^\ast$ hängt \emph{ausschließlich} vom Kostenverhältnis $\beta/\alpha$
	(Round-Trip- vs.\ Knotenkosten) und vom Verhältnis $d/p$ (Verzweigungsgrad vs.\
	Fortsetzungswahrscheinlichkeit) ab, \emph{nicht} jedoch von einer prinzipiellen
	REST-Inkompatibilität (Satz~\ref{satz:hst-zustandslosigkeit}). Für hochlatente
	Verbindungen (großes $\beta$, z.\,B.\ mobile Clients) oder Prozesse mit hoher
	Fortsetzungswahrscheinlichkeit $p$ (lange, vorhersehbare Sequenzen, etwa
	Sammelzustände gemäß Bemerkung~\ref{bem:hst-permutierbar}) wächst $h^\ast$.
	Abbildung~\ref{fig:optimaler-horizont} (rechts) zeigt diese Abhängigkeit als
	Stufenfunktion, da $h^\ast$ ganzzahlig ist. Insbesondere ist für Sammelzustände mit
	$r$ permutierbaren Teilschritten der Horizont $h^\ast\ge r$ (vollständige
	Dezentralisierung gemäß Bemerkung~\ref{bem:hst-permutierbar}) gemäß (c) genau dann
	optimal, wenn $\beta/\alpha\ge(d/p)^{r}$ -- für latenzkritische Anwendungen (großes
	$\beta/\alpha$) also bereits für moderate $r$ erfüllt, während für
	latenzunkritische Anwendungen (kleines $\beta/\alpha$) ein kleinerer Horizont
	$h^\ast<r$ kosteneffizienter ist und die Sammelzustands-Konformität partiell
	serverseitig verbleibt.
\end{bemerkung}

\begin{figure}[ht]
	\centering
	\includegraphics[width=\textwidth]{plots/plot6_optimaler_horizont.pdf}
	\caption{Links: Gesamtkosten $C(h)$ (logarithmische Skala) gemäß
	Definition~\ref{def:kostenmodell-horizont} für $\alpha=1$, $\beta=50$, $d=2$ und
	Fortsetzungswahrscheinlichkeiten $p\in\{0.3,0.5,0.7\}$. Die Sterne markieren das nach
	Satz~\ref{satz:optimaler-horizont} eindeutige Minimum $h^\ast$: Für $p=0.3$ und
	$p=0.5$ liegt es bei $h^\ast=2$, für $p=0.7$ bei $h^\ast=3$ -- in Übereinstimmung
	mit der geschlossenen Form aus (c), denn für $p=0.7$ gilt
	$\log_{2/0.7}(50)\approx3.83$, also $h^\ast=\lceil3.83\rceil-1=3$. Rechts: optimaler
	Horizont $h^\ast$ als Funktion des Kostenverhältnisses $\beta/\alpha$
	(logarithmische Abszisse) für $d=2$ und dieselben $p$-Werte, als Stufenfunktion
	gemäß (c). Größere $\beta/\alpha$ (teurere Round-Trips relativ zur Payload-Größe,
	z.\,B.\ bei hoher Netzwerklatenz) und größere $p$ (Prozesse mit tendenziell langen
	Aufrufsequenzen) führen zu größerem $h^\ast$; die Stufen markieren die in (c)
	bewiesenen Schwellenwerte $\beta/\alpha=(d/p)^{h+1}$.}
	\label{fig:optimaler-horizont}
\end{figure}

\subsection{Versionierung und Cache-Kohärenz via Signaturen}
\label{sec:hst-cache}

\begin{definition}[Signatur des Horizont-Subgraphen]
	\label{def:signatur-horizont}
	Sei $G_s^h=(V_s^h,E_s^h)$ ein Horizont-Subgraph mit einer fest gewählten,
	kanonischen Aufzählung $V_s^h=\{v_0,\ldots,v_{m-1}\}$ (z.\,B.\ in BFS-Reihenfolge
	gemäß Definition~\ref{def:horizont-subgraph}, mit $v_0=s$) und sei
	$\mathrm{idx}:\Sigma\to\{0,\ldots,|\Sigma|-1\}$ eine feste Indizierung des
	Operationsalphabets. Die \emph{Signatur} von $G_s^h$ ist
	\[
		\mathrm{Sig}(G_s^h) := \sum_{(v_i,\sigma,v_j)\in E_s^h}
		\Bigl(\, i \;+\; j\cdot m \;+\; \mathrm{idx}(\sigma)\cdot m^2 \,\Bigr)\cdot
		\bigl(2|\Sigma|\,m^2\bigr)^{\,\mathrm{rang}(v_i,\sigma,v_j)},
	\]
	wobei $\mathrm{rang}(\cdot)$ eine feste Aufzählung der Kanten $0,1,2,\ldots$
	bezeichnet. Diese Konstruktion übertragt das Prinzip der Spalten-Signaturen
	$\sigma_j=\sum_iA_{ij}2^i+j\cdot2^n$ aus~\cite{epp2026subgraph} -- dort für
	Adjazenzmatrizen von Aufrufsequenzen definiert -- auf die (Knoten-, Kanten-,
	Beschriftungs-)Tripel von $G_s^h$: Jede Kante wird auf einen Wert im Bereich
	$[0,2|\Sigma|m^2)$ abgebildet und die Kanten werden, analog zur stellenwertigen
	Kodierung in~\cite{epp2026subgraph}, mit unterschiedlichen „Stellen“ (Potenzen von
	$2|\Sigma|m^2$) gewichtet, sodass $\mathrm{Sig}$ auf der Menge aller Graphen mit
	höchstens $m$ Knoten über dem Operationsalphabet $\Sigma$ injektiv ist (Übertragung
	des Eindeutigkeitslemmas aus~\cite{epp2026subgraph} auf die hier verwendete
	stellenwertige Kodierung).
\end{definition}

\begin{bemerkung}[Cache-Kohärenz für Navigationsgraphen]
	\label{bem:cache-kohaerenz}
	Ein Client speichert $(G_s^h,\mathrm{Sig}(G_s^h))$ lokal. Wird der Prozessautomat
	$P$ -- etwa im Rahmen einer fachlichen Prozessänderung -- zu $P'$ aktualisiert,
	sodass sich $\delta$ in einem für $s$ relevanten Bereich ändert, so unterscheidet
	sich der neue Horizont-Subgraph $(G')_s^h$ von $G_s^h$ in mindestens einer Kante.
	Da $\mathrm{Sig}$ nach Definition~\ref{def:signatur-horizont} injektiv ist, folgt
	$\mathrm{Sig}((G')_s^h)\ne\mathrm{Sig}(G_s^h)$. Ein HST-Server kann die aktuelle
	Signatur in einem \texttt{ETag}-Header der HST-Repräsentation übermitteln; ein
	Client kann damit in $O(1)$ (Vergleich zweier Zahlen) erkennen, ob sein lokal
	zwischengespeicherter Horizont-Subgraph noch gültig ist, und nur im Falle einer
	Änderung eine vollständige Neuübertragung von $G_s^h$ anfordern
	(\texttt{If-None-Match}-Mechanismus gemäß~\cite{rfc7231}). Dies überträgt das in
	Abschnitt~\ref{sec:pattern} (Unterabschnitt „Zustandslosigkeit und SC-REST“)
	beschriebene Prinzip der \texttt{ETag}-basierten Versionsprüfung von einzelnen
	Ressourcenzuständen auf den \emph{gesamten, im Horizont enthaltenen
	Navigationsgraphen} -- ein Client kann damit nicht nur erkennen, ob sich
	\emph{seine} Ressource geändert hat, sondern auch, ob sich die für ihn relevanten
	\emph{Spielregeln} ($\delta$ eingeschränkt auf $G_s^h$) geändert haben.
\end{bemerkung}

\subsection{HST als Architekturstil: Generalisierung von REST}
\label{sec:hst-architekturstil}

Fielding definiert Architekturstile in~\cite{fielding2000} durch sukzessives Hinzufügen
von Beschränkungen (\emph{constraints}) zu einem „Nullstil“ ohne jede Beschränkung; REST
entsteht durch Hinzufügen von Client-Server, Zustandslosigkeit, Cache, Uniform
Interface, Layered System und (optional) Code-on-Demand. Die in
Abschnitt~\ref{sec:hst-vollstaendigkeit} bewiesene lokale Vollständigkeit motiviert eine
\emph{zusätzliche} Beschränkung, mit der sich HST formal als echte Verfeinerung von REST
einordnen lässt.

\begin{definition}[HST-Architekturstil]
	\label{def:hst-architekturstil}
	Der \emph{HST-Architekturstil} ist der REST-Architekturstil~\cite{fielding2000},
	erweitert um die folgende zusätzliche Beschränkung:
	\begin{quote}
		\textbf{Horizont-Lokalität.} Es existiert ein anwendungsweit oder
		ressourcenspezifisch konfigurierter Horizont $h\ge1$, sodass jede
		Repräsentation einer zustandsbehafteten Ressource zusätzlich zu den unmittelbar
		gültigen Übergängen (Uniform Interface / HATEOAS) den vollständigen
		$h$-Schritt-Horizont-Subgraphen $G_s^h$ gemäß
		Definition~\ref{def:horizont-subgraph} enthält, sodass jeder Client für
		Aufrufsequenzen der Länge $\le h$ lokale Vollständigkeit im Sinne von
		Satz~\ref{satz:lokale-vollstaendigkeit} besitzt.
	\end{quote}
\end{definition}

\begin{table}[ht]
	\centering
	\renewcommand{\arraystretch}{1.35}
	\caption{Die REST-Beschränkungen nach~\cite{fielding2000} und ihre Realisierung in
	SC-REST ($h=1$) bzw.\ HST ($h\ge1$ allgemein).}
	\label{tab:rest-vergleich}
	\begin{tabular}{@{}p{3.3cm}p{5.55cm}p{5.55cm}@{}}
		\toprule
		REST-Beschränkung & SC-REST ($h=1$) & HST ($h\ge1$, allgemein) \\
		\midrule
		Client-Server & unverändert & unverändert \\
		Zustands\-losigkeit &
		$\pi(id)=s$ als Ressourcenzustand (Abschn.~\ref{sec:pattern}) &
		unverändert gültig für beliebiges $h$ (Satz~\ref{satz:hst-zustandslosigkeit}) \\
		Cache &
		Standard-HTTP-Caching pro Ressourcenzustand &
		zusätzlich signaturbasierte Kohärenz für $G_s^h$
		(Abschn.~\ref{sec:hst-cache}) \\
		Uniform Interface (\mbox{HATEOAS}) &
		Links auf $\Sigma(\delta(s,\sigma))$, Horizont $1$ &
		vollständiger Horizont-Subgraph $G_s^h$, Horizont $h\ge1$ konfigurierbar \\
		Layered System & unverändert & unverändert \\
		Code-on-Demand & nicht betrachtet & nicht betrachtet \\
		\emph{(neu)}\newline Horizont-Lokalität &
		trivial erfüllt für Sequenzlänge $\le1$ &
		erfüllt für Sequenzlänge $\le h$ (Satz~\ref{satz:lokale-vollstaendigkeit}) \\
		\bottomrule
	\end{tabular}
\end{table}

\begin{bemerkung}
	Die Einordnung von HST als \emph{Erweiterung}, nicht als Ersetzung, von REST
	entspricht Fieldings Methodik~\cite{fielding2000}, Architekturstile inkrementell
	durch Hinzufügen von Beschränkungen zu definieren. Jede HST-konforme Architektur ist
	insbesondere REST-konform -- Satz~\ref{satz:hst-zustandslosigkeit} zeigt dies
	explizit für die Zustandslosigkeits-Beschränkung, die unter allen übrigen
	Beschränkungen invariant bleibt. HST ist somit, ebenso wie SC-REST, kein
	alternativer Stil zu REST, sondern eine \emph{Verfeinerung}: SC-REST verfeinert REST
	um Sequenzkonformität (Abschnitt~\ref{sec:pattern}), und HST verfeinert SC-REST
	zusätzlich um Horizont-Lokalität. Die Beziehung
	\[
		\text{REST} \;\supseteq\; \text{SC-REST} \;\supseteq\; \text{HST}
	\]
	(als Mengen von Architekturen, die die jeweiligen Beschränkungen erfüllen) ist
	gemäß Bemerkung~\ref{bem:sc-rest-spezialfall} sogar eine Gleichheit zwischen
	SC-REST und HST für $h=1$: SC-REST ist also nicht nur \emph{enthalten} in HST,
	sondern der Grenzfall $h=1$ \emph{ist} HST.
\end{bemerkung}

\subsection{Referenzimplementierung: Berechnung des Horizont-Subgraphen}
\label{sec:hst-implementierung}

Listing~\ref{lst:hst-bfs} erweitert die Middleware aus
Listing~\ref{lst:transition}/\ref{lst:flask} um eine Breitensuche, die den
Horizont-Subgraphen $G_s^h$ gemäß Definition~\ref{def:horizont-subgraph} berechnet, und
um eine Funktion, die diesen Subgraphen gemäß Definition~\ref{def:hst-repraesentation} in
die HST-Repräsentation einbettet.

\begin{lstlisting}[language=Python, label=lst:hst-bfs, caption=Breitensuchbasierte Berechnung des Horizont-Subgraphen $G_s^h$ und Einbettung in die HST-Repraesentation]
def horizon_subgraph(process: ProcessAutomaton,
                      s: str, h: int):
    """
    Berechnet G_s^h = (V, E) gemaess Definition
    "h-Schritt-Horizont-Subgraph" per Breitensuche.

    V: Menge erreichbarer Zustaende (Knoten)
    E: Liste von (s1, operation, s2)-Tripeln (Kanten)
    """
    V = {s}
    E = []
    frontier = {s}

    for _ in range(h):
        next_frontier = set()
        for s1 in frontier:
            for op in process.allowed_operations(s1):
                s2 = process.delta[(s1, op)]
                E.append((s1, op, s2))
                if s2 not in V:
                    V.add(s2)
                    next_frontier.add(s2)
        frontier = next_frontier
        if not frontier:
            break  # Horizont reicht bis F oder Sackgasse

    return V, E


def sst_representation(process: ProcessAutomaton,
                        instance_id: str, h: int = 2):
    """
    Erstellt die HST-Repraesentation R(id, s, h) gemaess
    Definition "HST-Repraesentation".
    """
    s = process.pi[instance_id]
    V, E = horizon_subgraph(process, s, h)

    return {
        "orderId": instance_id,
        "state": s,
        "_horizon": h,
        "_graph": {
            "nodes": sorted(V),
            "edges": [
                {"from": s1, "op": op, "to": s2}
                for (s1, op, s2) in E
            ],
        },
        # h=1-Teilmenge entspricht den bisherigen
        # HATEOAS-Links (Listing 4, Bemerkung "SC-REST
        # ist Spezialfall h=1 von HST")
        "_links": hateoas_links(process, s, instance_id),
    }
\end{lstlisting}

\begin{bemerkung}
	Die Laufzeit von \texttt{horizon\_subgraph} ist $O(|V_s^h|\cdot d)= O(d^{h+1})$
	(Lemma~\ref{lem:horizont-groesse}), da jede der $h$ BFS-Schichten höchstens
	$|V_s^h|$ Knoten mit jeweils höchstens $d$ ausgehenden Operationen betrachtet; bei
	Hashtabellen-basierter Implementierung von $\delta$ (Satz~\ref{satz:laufzeit-pruefung})
	ist jeder Zugriff $\delta[(s_1,\sigma)]$ in erwarteter konstanter Zeit möglich. Für
	die in Abbildung~\ref{fig:horizont-groesse} betrachteten praxisrelevanten Horizonte
	$h\le4$ und $d\le3$ ist dieser Aufwand vernachlässigbar gegenüber der ohnehin
	anfallenden Fachlogik-Ausführung in Definition~\ref{def:hst-server}.
\end{bemerkung}

\section{Zusammenfassung}

Diese Arbeit hat mit SC-REST ein formales Pattern für sequenzgebundene REST-Schnittstellen
entwickelt (Abschnitte~\ref{sec:pattern}--\ref{sec:implementierung}) und dieses Pattern in
Abschnitt~\ref{sec:sst} zu \emph{HST (Horizon State Transfer)} weiterentwickelt: SC-REST
übermittelt pro Antwort die im aktuellen Zustand zulässigen Folgeoperationen ($1$-Schritt-
Horizont), während HST den vollständigen $h$-Schritt-Horizont-Subgraphen $G_s^h$
überträgt und damit lokale Vollständigkeit für Aufrufsequenzen der Länge $\le h$
herstellt (Satz~\ref{satz:lokale-vollstaendigkeit}), ohne die Zustandslosigkeit nach
Fielding~\cite{fielding2000} zu beeinträchtigen (Satz~\ref{satz:hst-zustandslosigkeit}).
SC-REST ist dabei formal der Spezialfall $h=1$ von HST
(Bemerkung~\ref{bem:sc-rest-spezialfall}); beide Pattern teilen daher dieselbe
Anwendungsgrundlage, die im Folgenden zusammengefasst wird, bevor in
Abschnitt~\ref{sec:ausblick} -- nun unter besonderer Berücksichtigung von HST --
Erweiterungen und Forschungsrichtungen diskutiert werden.

\subsection{Anwendungen}

Das SC-REST-Pattern (und, mit größerem Horizont, HST) eignet sich insbesondere für:
\begin{itemize}
	\item E-Commerce-Checkout-Prozesse (Warenkorb, Versand, Zahlung, Bestätigung) wie in
	      Beispiel~\ref{bsp:bestellprozess},
	\item mehrstufige Antrags- und Genehmigungsprozesse in Verwaltung und Finanzwesen
	      (Einreichung, Prüfung, Freigabe, Auszahlung),
	\item Onboarding-Flows mit zwingender Reihenfolge (Identitätsprüfung vor
	      Vertragsabschluss vor Zahlungseinrichtung),
	\item zustandsbehaftete Geräteprotokolle, die über REST gekapselt werden
	      (Initialisierung, Kalibrierung, Messung, Abschluss),
	\item Multi-Step-Formulare und Wizards in Web-Anwendungen, bei denen
	      Schrittreihenfolgen serverseitig erzwungen werden müssen, um inkonsistente
	      Zwischenzustände zu verhindern.
\end{itemize}

Für HST (Abschnitt~\ref{sec:sst}) treten insbesondere solche Anwendungen hinzu, bei denen
\emph{Planungsvorsprung} fachlich wertvoll ist: Mobile Clients mit instabiler Verbindung
(großes $\beta$ in Definition~\ref{def:kostenmodell-horizont}, daher größeres $h^\ast$
nach Satz~\ref{satz:optimaler-horizont}), grafische Prozess-Editoren und Wizard-UIs, die
dem Nutzer \emph{vorab} alle erreichbaren Folgeschritte mehrerer Ebenen anzeigen möchten
(„Fortschrittsanzeige“ als direkte UI-Repräsentation von $G_s^h$), sowie -- gemäß
Bemerkung~\ref{bem:hst-permutierbar} -- Sammelzustände mit permutierbaren Teilschritten,
für die HST mit $h\ge r$ die in Abschnitt~\ref{sec:permutierbar} entwickelte
signaturbasierte Konformitätsprüfung vollständig auf den Client verlagert.

\subsection{Ausblick}
\label{sec:ausblick}

Das in dieser Arbeit vorgestellte SC-REST-Pattern und seine Verallgemeinerung zu HST
(Abschnitt~\ref{sec:sst}) bilden die formale Grundlage für zahlreiche Erweiterungen und
Forschungsrichtungen, die im Folgenden ausführlich diskutiert werden. Die ersten zwölf
Unterabschnitte betreffen primär SC-REST und seine in Abschnitt~\ref{sec:analyse} bis
\ref{sec:implementierung} entwickelten Resultate; die abschließenden Unterabschnitte
widmen sich speziell den durch HST eröffneten Fragestellungen.

\subsubsection{Formale Spezifikationssprache für Prozessautomaten}

Die Definitionen~\ref{def:prozessautomat}--\ref{def:konformitaet} legen nahe, $\delta$
nicht imperativ (wie in Listing~\ref{lst:flask}) zu kodieren, sondern \emph{deklarativ} in
einer maschinenlesbaren Spezifikation. Ein naheliegender Ansatz ist eine Erweiterung von
OpenAPI um ein Feld \texttt{x-process-state}, das pro Operation den Vorbedingungszustand
(\texttt{requires-state}) und den Folgezustand (\texttt{transitions-to}) angibt. Aus einer
solchen Spezifikation ließen sich automatisch (1) die Hashtabelle $\delta$ für
Satz~\ref{satz:laufzeit-pruefung}, (2) Client-SDKs mit typsicheren Zustandsübergängen, (3)
interaktive API-Dokumentation mit Zustandsdiagrammen analog zu
Abbildung~\ref{fig:automat}, und (4) Testsuiten generieren, die systematisch alle
$\delta(s,\sigma)$ sowie alle erwarteten \texttt{409}-Fälle (alle $(s,\sigma)\notin
\mathrm{dom}(\delta)$) abdecken. Eine solche DSL würde die in
Bemerkung~\ref{bem:prefetch-kosten} geforderte Begrenzung $d(s)\leq D_{\max}$ bereits zur
Designzeit statisch prüfbar machen.

\subsubsection{Verteilte Prozessinstanzen und das Saga-Pattern}

Reale Geschäftsprozesse erstrecken sich häufig über mehrere Microservices, von denen jeder
nur einen Teil von $\Sigma$ implementiert. Eine natürliche Erweiterung besteht darin, den
globalen Prozessautomaten $P$ als \emph{Choreographie} mehrerer lokaler Automaten
$P_1,\ldots,P_n$ zu modellieren, deren Zustände $S_i$ über Ereignisse synchronisiert
werden -- eine Verbindung zum klassischen Saga-Pattern für verteilte Transaktionen
\cite{garcia1987}. Im Unterschied zu klassischen Sagas, die primär
\emph{Kompensationsaktionen} bei Fehlern definieren, würde eine SC-REST-Choreographie
zusätzlich \emph{präventiv} verhindern, dass ein Client überhaupt eine Operation auslöst,
deren Vorbedingungen in einem anderen Service noch nicht erfüllt sind. Hierzu wäre
Satz~\ref{satz:korrektheit} auf ein verteiltes Setting zu verallgemeinern, in dem $\pi(id)$
nicht von einem einzelnen Server, sondern von einer verteilten, eventuell replizierten
Datenstruktur verwaltet wird -- mit den entsprechenden Konsistenzfragen (siehe Punkt
„Multi-Tenant“ weiter unten).

\subsubsection{Nichtdeterministische Automaten und parallele Teilprozesse}

Definition~\ref{def:prozessautomat} fordert, dass $\delta$ eine \emph{Funktion} ist
(deterministisch). Viele Prozesse enthalten jedoch \emph{parallele} Zweige, die
unabhängig voneinander und potenziell gleichzeitig fortschreiten (im Unterschied zu den in
Abschnitt~\ref{sec:permutierbar} behandelten \emph{sequentiell, aber in beliebiger
Reihenfolge} auszuführenden Teilschritten). Eine natürliche Erweiterung ist der Übergang
von Prozessautomaten zu \emph{Petri-Netzen}~\cite{murata1989}: Stellen (places)
modellieren Teilprozesse, Marken (tokens) den Fortschritt, und Transitionen werden erst
„feuerbar“, wenn alle Vorbedingungs-Stellen markiert sind. Die Konformitätsprüfung
(Satz~\ref{satz:pfadcharakterisierung}) würde dann zu einer Erreichbarkeitsprüfung im
Markierungsgraphen des Petri-Netzes, deren Komplexität im Allgemeinen deutlich höher ist
als $\Theta(k)$ -- ein Forschungsfeld, in dem die Ãœbertragung der
Subgraph Algorithmus-Technik aus Satz~\ref{satz:permutierbar-subgraph} auf
Markierungsvektoren ein vielversprechender Ansatzpunkt für zukünftige Arbeiten ist.

\subsubsection{Sicherheit: Geschützte Zustandstoken}

In der hier vorgestellten Form wird $\pi(id)$ serverseitig gespeichert. In hochskalierbaren
oder Edge-Architekturen kann es vorteilhaft sein, den Zustand $s$ direkt im Client zu
halten -- kodiert als signiertes JSON Web Token (JWT), das bei jeder Anfrage mitgesendet
wird. Dies würde die Statelessness im klassischen Sinn vollständig wiederherstellen
(\emph{kein} serverseitiger Zustand pro Instanz), erfordert jedoch eine kryptografische
Signatur, die Manipulation des Zustands durch den Client ausschließt (sonst könnte ein
Client durch direktes Setzen von $s:=s_4$ die Prüfung $\delta(s,\sigma)$ umgehen). Die
formale Korrektheitsaussage aus Satz~\ref{satz:korrektheit} bliebe in diesem Modell
gültig, sofern die Signaturprüfung als zusätzliche, der eigentlichen
$\delta$-Prüfung vorgeschaltete Operation modelliert wird, deren Fehlschlagen ebenfalls zu
einer (jedoch von \texttt{sequence-violation} zu unterscheidenden,
\texttt{401}/\texttt{403}-artigen) Fehlerantwort führt.

\subsubsection{Adaptive Vorhersage mittels maschinellem Lernen}

Das Latenzmodell aus Definition~\ref{def:latenzmodell} und Satz~\ref{satz:latenzgewinn}
geht von einem festen, durch $N(s)$ gegebenen Kandidatensatz für das Prefetching aus. Bei
hohem Verzweigungsgrad $d(s)$ (Bemerkung~\ref{bem:prefetch-kosten}) ist es jedoch
ineffizient, \emph{alle} $|N(s)|$ Kandidaten gleich zu priorisieren. Stattdessen könnte ein
Markov-Modell höherer Ordnung -- basierend auf der Historie vorheriger Prozessinstanzen --
eine Wahrscheinlichkeitsverteilung $P(\sigma \mid s, \text{Kontext})$ über die nächste
Operation schätzen und das Prefetching nur für die $w$ wahrscheinlichsten Kandidaten
durchführen (\emph{Top-$w$-Prefetching}). Satz~\ref{satz:latenzgewinn} ließe sich dann zu
einem \emph{erwarteten} Latenzgewinn
$E[\Delta] = \sum_{\sigma\in\Sigma(s)} P(\sigma\mid s)\cdot
\min(T_{\mathrm{think}},T_{\mathrm{prep}}(\delta(s,\sigma)))$ verfeinern, und die Wahl von
$w$ würde einen expliziten Kompromiss zwischen erwartetem Latenzgewinn und
Ressourcenverbrauch (Bemerkung~\ref{bem:prefetch-kosten}) parametrisieren. Die in
Abbildung~\ref{fig:lernkurve} gezeigte Lernkurve liefert hierfür bereits die empirische
Grundlage: Die mit zunehmender Prozessreife sinkende Fehlerrate $p(n)$ entspricht einer
zunehmend „spitzeren“ Verteilung $P(\sigma\mid s)$.

\subsubsection{Verifikation und Model Checking}

Da $P=(S,\Sigma,\delta,s_0,F)$ ein endlicher Automat ist, lassen sich Eigenschaften wie
„jeder Zustand $s\in S\setminus F$ besitzt mindestens einen Pfad nach $F$“
(\emph{Lebendigkeit}, keine „Sackgassen“) oder „$F$ ist von $s_0$ aus stets erreichbar“
mit Standardverfahren der Erreichbarkeitsanalyse in $O(|S|+|E|)$ prüfen (Tiefen- bzw.
Breitensuche auf $G_P$, Definition~\ref{def:prozessgraph}). Für die in
Abschnitt~\ref{sec:permutierbar} diskutierten Sammelzustände mit ihrer impliziten
$2^r$-Struktur (Lemma~\ref{lem:zustandsexplosion}) wären jedoch symbolische Verfahren
(BDD-basiertes Model Checking, LTL/CTL-Spezifikationen) notwendig, um Eigenschaften wie
„jede Reihenfolge der $r$ optionalen Teilschritte führt zu einem konsistenten
Folgezustand“ ohne explizite Aufzählung aller $2^r$ Zwischenzustände zu verifizieren --
ein direkter Anknüpfungspunkt für die formale Methode dieser Arbeit.

\subsubsection{Integration mit GraphQL, gRPC und AsyncAPI}

Das SC-REST-Pattern wurde für REST/HTTP mit JSON-Repräsentationen formuliert. Die
zugrundeliegenden Definitionen~\ref{def:prozessautomat}--\ref{def:sc-rest-server} sind
jedoch protokollunabhängig: Für GraphQL ließe sich $\Sigma$ als Menge von Mutationen
interpretieren, wobei \texttt{409}-artige Fehler als typisierte GraphQL-Errors mit einem
analogen \texttt{sequence-violation}-Feld kodiert würden; für gRPC entsprächen die
Operationen RPC-Methoden, und Sequenzverletzungen würden über den
\texttt{FAILED\_PRECONDITION}-Statuscode signalisiert, dessen Metadaten die Menge
$\Sigma(s)$ trügen. Für ereignisgetriebene Architekturen (AsyncAPI) wäre $\Sigma$ die
Menge der konsumierten Nachrichtentypen, und eine Sequenzverletzung würde als
\emph{Dead-Letter}-Ereignis mit Verweis auf $\Sigma(s)$ veröffentlicht. In allen Fällen
blieben Satz~\ref{satz:korrektheit} (Korrektheit) und Satz~\ref{satz:laufzeit-pruefung}
(Laufzeit) unverändert gültig, da beide ausschließlich auf der abstrakten Struktur von
$P$, nicht auf HTTP-spezifischen Details, beruhen.

\subsubsection{Hierarchische Automaten und Subprozesse}

Komplexe Prozesse lassen sich häufig in Subprozesse zerlegen, die selbst wieder
Prozessautomaten gemäß Definition~\ref{def:prozessautomat} sind. Ein Zustand $s\in S$ des
übergeordneten Automaten könnte dann „intern“ einen vollständigen Subprozessautomaten
$P_s=(S_s,\Sigma_s,\delta_s,s_0^{(s)},F_s)$ repräsentieren, wobei der Übergang aus $s$ im
übergeordneten Automaten erst erfolgt, wenn $P_s$ einen Zustand in $F_s$ erreicht hat. Der
Prozessgraph $G_P$ (Definition~\ref{def:prozessgraph}) würde dadurch zu einem
\emph{hierarchischen Graphen}, in dem jeder „Superknoten“ selbst einen Graphen $G_{P_s}$
enthält. Die Konformitätsprüfung einer Aufrufsequenz würde rekursiv: Auf jeder
Hierarchieebene wird Satz~\ref{satz:pfadcharakterisierung} angewendet, und für
Sammelzustände mit permutierbaren Teilschritten (Abschnitt~\ref{sec:permutierbar}) auf
jeder Ebene unabhängig Satz~\ref{satz:permutierbar-subgraph}. Dies eröffnet die
Möglichkeit, den Subgraph Algorithmus aus~\cite{epp2026subgraph} \emph{rekursiv} auf
verschachtelte Subgraphen anzuwenden -- ein direkter struktureller Bezug zur in
\cite{epp2026subgraph} diskutierten AST-Analyse, bei der Teilbäume (Subgraphen)
ebenfalls rekursiv verglichen werden.

\subsubsection{Standardisierung: Vorschlag für eine OpenAPI-Erweiterung „x-sc-rest“}

Die in dieser Arbeit entwickelten Definitionen könnten Grundlage eines formalen
Standardisierungsvorschlags sein, etwa als OpenAPI-Erweiterungsfeld
\texttt{x-sc-rest} oder als eigenständiges Internet-Draft analog zu RFC~7807
\cite{rfc7807}. Ein solcher Standard müsste mindestens spezifizieren: (1) das Format zur
Deklaration von $S,\Sigma,\delta,s_0,F$ (vgl. „Formale Spezifikationssprache“ oben), (2)
den Problem-Detail-Typ \texttt{sequence-violation} inklusive des Feldes
\texttt{allowedOperations}, (3) die HATEOAS-Link-Relationstypen für Folgeoperationen, und
(4) optionale Erweiterungsfelder für Sammelzustände (Definition~\ref{def:sammelzustand}),
etwa \texttt{x-sc-rest-collection} mit den Feldern \texttt{operations} (entsprechend
$\Sigma_{\mathrm{opt}}$) und \texttt{completion} (entsprechend $\sigma_{\mathrm{done}}$).
Ein solcher Standard würde die in Pautasso et al.~\cite{pautasso2008} diskutierte Debatte
zwischen „RESTful“ und „Big Web Services“ (WS-BPEL) um eine dritte Position ergänzen:
prozessbewusste REST-APIs, die die Vorteile beider Welten -- die Einfachheit von REST und
die Prozessgarantien von WS-BPEL-Choreographien -- formal vereinen.

\subsubsection{Experimentelle Validierung in Produktionssystemen}

Die in Abschnitt~\ref{sec:empirisch} gezeigten Diagramme basieren auf simulierten Daten.
Eine naheliegende Folgearbeit ist die Instrumentierung einer SC-REST-Referenzimplementierung
(Abschnitt~\ref{sec:implementierung}) in einem produktiven Microservice-Umfeld
\cite{newman2021}, um (1) die in Satz~\ref{satz:latenzgewinn} postulierte
Latenzreduktion unter realer Lastverteilung von $T_{\mathrm{think}}$ und
$T_{\mathrm{prep}}$ zu vermessen, (2) die tatsächliche Verteilung der
Verzweigungsgrade $d(s)$ in realen Prozessmodellen zu erheben (als Eingabe für
Bemerkung~\ref{bem:prefetch-kosten}), und (3) die in Lemma~\ref{lem:zustandsexplosion}
postulierte Zustandsraum-Explosion für reale Sammelzustände (typische Werte für $r$) zu
quantifizieren. Eine offene Benchmark-Suite, die verschiedene Prozessautomaten
$P$ unterschiedlicher Größe $|S|$, $|\Sigma|$ und Verzweigungsgrade enthält, würde
zudem den systematischen Vergleich verschiedener Implementierungsstrategien für $\delta$
(Hashtabelle, Entscheidungsbaum, kompilierter endlicher Automat) ermöglichen.

\subsubsection{Energieeffizienz durch reduzierte Fehlversuche}

Jede \texttt{409}-Antwort gemäß Satz~\ref{satz:korrektheit}(c) repräsentiert eine
„vergebliche“ Anfrage: Netzwerkübertragung, TLS-Handshake (falls keine Connection
wiederverwendet wird), Server-seitige Verarbeitung der Anfrage und der Fehlerantwort --
alles ohne fachlichen Nutzen. Die in Abbildung~\ref{fig:lernkurve} gezeigte sinkende
Fehlerrate $p(n)$ impliziert daher nicht nur einen Latenz-, sondern auch einen
Energieeffizienzgewinn: Pro vermiedener \texttt{409}-Antwort wird die gesamte
Verarbeitungskette einmal weniger durchlaufen. Eine quantitative Abschätzung dieses
Effekts -- etwa in Joule pro vermiedener Anfrage, aggregiert über Rechenzentren mit
hohem API-Durchsatz -- wäre ein Beitrag zur Diskussion um „Green Software Engineering“ und
ergänzt die in Bemerkung~\ref{bem:prefetch-kosten} diskutierte Kostenseite des
Prefetchings: Während Prefetching gezielt \emph{zusätzliche} Berechnungen vorzieht, um
Latenz zu sparen, spart die Sequenzkontrolle selbst Berechnungen \emph{vollständig} ein,
indem sie von vornherein unmögliche Anfragen früh und billig (Satz
\ref{satz:laufzeit-pruefung}, $O(1)$) abweist.

\subsubsection{Multi-Tenant-Architekturen und verteilte Zustandsspeicher}

Definition~\ref{def:prozessinstanz} setzt voraus, dass $\pi: \text{ID}\rightharpoonup S$
konsistent gelesen und geschrieben werden kann. In horizontal skalierten Systemen wird
$\pi$ typischerweise in einem verteilten Schlüssel-Wert-Speicher (z.\,B. Redis-Cluster)
gehalten. Damit stellt sich die Frage, welche Konsistenzgarantien für $\pi$ erforderlich
sind, damit Satz~\ref{satz:korrektheit} weiterhin gilt: Wird $\pi(id)$ unter
Eventual Consistency repliziert, könnten zwei nahezu gleichzeitige Anfragen für dieselbe
Prozessinstanz $id$ auf unterschiedlichen Replikaten denselben Ausgangszustand $s$ sehen
und beide einen (jeweils für sich konformen) Übergang $\delta(s,\sigma)$ und
$\delta(s,\sigma')$ ausführen -- ein klassisches \emph{Lost-Update}-Problem, das die
Eindeutigkeit des Laufs aus Definition~\ref{def:lauf} verletzt. Eine Erweiterung dieser
Arbeit müsste daher entweder (a) starke Konsistenz für $\pi(id)$ fordern (z.\,B. mittels
eines verteilten Locks oder einer Single-Writer-Partitionierung nach $id$, was angesichts
der CAP-Theorem-Implikationen Verfügbarkeitseinbußen bei Partitionierung bedeutet), oder
(b) Satz~\ref{satz:korrektheit} zu einer \emph{optimistischen} Variante verfeinern, bei
der konkurrierende Übergänge mittels eines Versionszählers (analog zu \texttt{ETag}/
\texttt{If-Match}) erkannt und einer der beiden Übergänge nachträglich mit
\texttt{409 Conflict} abgewiesen wird -- eine Idempotenz- und
Konfliktbehandlungsstrategie, die eng mit der in Abschnitt~\ref{sec:robustheit}
diskutierten Replay-Sicherheit (Satz~\ref{satz:replay}) zusammenhängt und in
\cite{helland2007} für verteilte Transaktionen allgemein diskutiert wird.

\subsubsection{Adaptiver Horizont: Lernen von $h^\ast$ pro Client und Prozesstyp}

Satz~\ref{satz:optimaler-horizont} liefert $h^\ast$ als Funktion der Parameter
$\alpha,\beta,d,p$ -- in der Praxis sind insbesondere $\beta$ (Round-Trip-Latenz) und $p$
(Fortsetzungswahrscheinlichkeit) jedoch nicht global konstant, sondern hängen vom
\emph{Client} (mobile Verbindung vs.\ Rechenzentrums-zu-Rechenzentrums-Aufruf, also
unterschiedliches $\beta$) und vom \emph{Prozesstyp} bzw.\ sogar vom \emph{individuellen
Nutzer} (manche Nutzer brechen Bestellvorgänge häufig nach dem ersten Schritt ab, andere
durchlaufen sie zuverlässig bis zum Ende, also unterschiedliches $p$) ab. Dies setzt die
in Abschnitt~\ref{sec:ausblick} (Unterabschnitt „Adaptive Vorhersage mittels
maschinellem Lernen“) für $T_{\mathrm{think}}$ und $T_{\mathrm{prep}}$ skizzierte Idee
fort: Ein lernendes System könnte $\hat\beta$ aus den \texttt{RTT}-Werten der
TCP-Verbindung des jeweiligen Clients und $\hat p$ aus der historischen
Sequenzlängenverteilung des jeweiligen Nutzers (oder, bei neuen Nutzern, aus einem
Prozesstyp-Prior) schätzen und daraus gemäß der geschlossenen Form in
Satz~\ref{satz:optimaler-horizont}(c) einen \emph{personalisierten} Horizont
$\hat h^\ast(\text{client},\text{user})$ berechnen, der pro Anfrage als
\texttt{?horizon=}-Parameter (Definition~\ref{def:hst-server}) gesetzt wird. Im
Gegensatz zu den meisten ML-Einsatzfeldern in zustandsbehafteten APIs ist diese Anwendung
\emph{folgenlos im Fehlerfall}: Eine falsche Schätzung von $\hat h^\ast$ führt nach
Satz~\ref{satz:hst-zustandslosigkeit} höchstens zu suboptimalen Kosten $C(\hat
h^\ast)>C(h^\ast)$, niemals zu einer inkorrekten oder inkonsistenten Antwort -- das
Lernproblem ist also ein reines Optimierungs-, kein Korrektheitsproblem, was den Einsatz
von ML hier mit deutlich geringerem Risiko behaftet als etwa bei der in
\cite{murata1989}-artigen Systemen üblichen Zustandsvorhersage.

\subsubsection{Informationsfreigabe und Zugriffskontrolle im Horizont-Subgraphen}

Definition~\ref{def:hst-repraesentation} setzt implizit voraus, dass \emph{der gesamte}
Horizont-Subgraph $G_s^h$ für den anfragenden Client offenlegbar ist. In Prozessen mit
rollenbasierter Autorisierung -- etwa Genehmigungs-Workflows, in denen ein
Sachbearbeiter nur die für seine Rolle bestimmten Übergänge sehen darf, während
Eskalationspfade (z.\,B.\ \texttt{POST /reject-and-escalate}) nur für Vorgesetzte
sichtbar sein sollen -- würde eine ungefilterte Übertragung von $G_s^h$ dem Client
\emph{Wissen über die Struktur des Prozessautomaten} offenlegen, das über seine eigenen
Handlungsmöglichkeiten hinausgeht: Allein die \emph{Existenz} eines
Eskalationszustands im Horizont könnte ein Informationsleck darstellen (etwa wenn die
bloße Existenz eines „Sonderkonditionen“-Zustands für reguläre Kunden vertraulich sein
soll). Eine konsequente Erweiterung von Definition~\ref{def:hst-server} wäre daher eine
\emph{rollenfilterte} Variante
\[
	G_s^h\!\restriction_\rho := \bigl(V_s^h\cap\mathrm{sichtbar}(\rho),\;
	E_s^h\cap(\mathrm{sichtbar}(\rho)\times\Sigma\times\mathrm{sichtbar}(\rho))\bigr),
\]
wobei $\rho$ die Rolle des Clients und $\mathrm{sichtbar}(\rho)\subseteq S$ die für $\rho$
sichtbaren Zustände bezeichnet. Eine formale Folgearbeit müsste klären, ob
$G_s^h\!\restriction_\rho$ weiterhin die in Satz~\ref{satz:lokale-vollstaendigkeit}
gezeigte lokale Vollständigkeit \emph{für die dem Client tatsächlich erlaubten
Operationen} besitzt -- vermutlich ja, sofern $\mathrm{sichtbar}(\rho)$ unter $\delta$
für die Operationen aus $\Sigma_\rho\subseteq\Sigma$ abgeschlossen ist -- und wie sich
dies auf die in Definition~\ref{def:signatur-horizont} definierte Signatur auswirkt
(vermutlich müsste $\mathrm{Sig}(G_s^h\!\restriction_\rho)$ rollenabhängig berechnet
werden, was die Cache-Kohärenz aus Abschnitt~\ref{sec:hst-cache} um eine zusätzliche
Dimension -- die Rolle -- erweitert).

\subsubsection{Komposition zu einem globalen Anwendungsgraphen}

Reale Anwendungen bestehen aus mehreren Prozessautomaten $P_1,\ldots,P_n$ -- etwa einem
Bestellprozess (Beispiel~\ref{bsp:bestellprozess}), einem Retourenprozess und einem
Reklamationsprozess, die über gemeinsame Ressourcen (dieselbe Bestellung) verknüpft sind.
Eine natürliche Erweiterung von Definition~\ref{def:horizont-subgraph} wäre ein
\emph{globaler Anwendungsgraph} $G_{\mathrm{App}} := \bigsqcup_{i=1}^n P_i \cup
E_{\mathrm{cross}}$, wobei $E_{\mathrm{cross}}$ \emph{Querverweis-Kanten} zwischen
Zuständen verschiedener Prozessautomaten enthält (z.\,B.\ vom Endzustand $s_5\in
F_{\text{Bestellung}}$ zum Startzustand $s_0'\in S_{\text{Retoure}}$, gemäß der Operation
\texttt{POST /orders/\{id\}/returns}). Der Horizont-Subgraph $G_s^h$ eines Zustands $s$
würde dann gegebenenfalls \emph{Zustände anderer Prozessautomaten} als Randknoten
enthalten, sobald $h$ groß genug ist, um $E_{\mathrm{cross}}$-Kanten zu erreichen. Dies
würde HATEOAS von der Navigation \emph{innerhalb} eines Prozesses zur Navigation
\emph{zwischen} Prozessen verallgemeinern -- formal müsste hierzu
Lemma~\ref{lem:horizont-groesse} auf $G_{\mathrm{App}}$ mit einem
\emph{kompositionalen} Verzweigungsgrad $d_{\mathrm{App}}=\max(d_1,\ldots,d_n,
|E_{\mathrm{cross}}|)$ verallgemeinert werden, und Satz~\ref{satz:optimaler-horizont}
müsste klären, ob ein einheitlicher Horizont $h^\ast$ über Prozessgrenzen hinweg
sinnvoll ist oder ob $E_{\mathrm{cross}}$-Kanten mit einem eigenen, typischerweise
größeren Kostenfaktor $\alpha_{\mathrm{cross}}>\alpha$ (Cross-Service-Aufruf statt
lokaler Tabellenzugriff) zu gewichten sind.

\subsubsection{Graphkompression für große Horizonte mittels Signaturen}

Für $h\ge r$ (Bemerkung~\ref{bem:hst-permutierbar}) enthält $G_{s_c}^h$ für einen
Sammelzustand $s_c$ den vollständigen $2^r$-Knoten-Entscheidungsbaum aller Teilmengen
$T\subseteq\{1,\ldots,r\}$ -- nach Lemma~\ref{lem:zustandsexplosion} potenziell sehr
groß. Viele dieser Teilbäume sind jedoch \emph{isomorph}: Der von einem Zwischenzustand
$s_T$ aus erreichbare Teilgraph hängt (bei symmetrischen optionalen Teilschritten) nur
von $|\{1,\ldots,r\}\setminus T|$, der Anzahl noch ausstehender Teilschritte, ab, nicht
von $T$ selbst. Eine signaturbasierte Kompression könnte daher isomorphe Teilbäume durch
\emph{Verweise} auf eine einzige kanonische Repräsentation ersetzen, identifiziert über
die Signatur $\mathrm{Sig}(\cdot)$ aus Definition~\ref{def:signatur-horizont}: Statt
$\binom{r}{k}$ strukturell identischer Teilbäume für $|T|=k$ würde nur \emph{ein}
Teilbaum übertragen und $\binom{r}{k}$ Referenzen darauf. Dies entspricht dem
klassischen Prinzip der \emph{strukturellen Teilung} (\emph{structure sharing}) in
funktionalen Datenstrukturen, übertragen auf die HST-Repräsentation; eine formale
Analyse müsste zeigen, dass diese Kompression die in
Bemerkung~\ref{bem:horizont-groesse-praxis} festgestellte exponentielle
Übertragungsgröße $|V_s^h|=O(d^h)$ auf eine \emph{polynomielle} Größe der komprimierten
Repräsentation reduziert -- ein direktes Analogon zur in
Bemerkung~\ref{bem:permutierbar-fazit} gezeigten Raumersparnis von $O(2^r)$ auf $O(r)$
für die serverseitige Kodierung, nun angewandt auf die clientseitig übertragene
HST-Repräsentation.

\subsubsection{Beziehung zu GraphQL: clientseitig verhandelter Horizont}

Der in Abschnitt~\ref{sec:ausblick} (Unterabschnitt „Integration mit GraphQL, gRPC und
AsyncAPI“) angesprochene Vergleich mit GraphQL gewinnt durch HST eine neue Dimension:
GraphQL erlaubt dem Client, die \emph{Tiefe} verschachtelter Relationen einer Anfrage
selbst zu spezifizieren (etwa „lade einen Autor, dessen Bücher, und für jedes Buch dessen
Rezensionen“ -- eine Tiefe von $2$). Der Horizont $h$ in
Definition~\ref{def:hst-repraesentation} ist strukturell analog: Ein HST-Server könnte
$h$ nicht als statische Konfiguration, sondern als \emph{vom Client pro Anfrage
spezifizierten Parameter} behandeln (wie bereits in Satz~\ref{satz:hst-zustandslosigkeit}
unter Punkt (3) als Option erwähnt), wodurch der Client den in
Satz~\ref{satz:optimaler-horizont} berechneten Tradeoff selbst -- abhängig von seiner
\emph{eigenen} Einschätzung von $\beta$ (seiner Verbindungsqualität) und $p$ (seinem
geplanten Workflow) -- vornehmen könnte, statt sich auf eine serverseitige Schätzung von
$\hat\beta,\hat p$ zu verlassen. Eine GraphQL-artige Anfrage könnte dann etwa lauten:
\texttt{GET /orders/\{id\}?horizon=3\&fields=state,total}, wobei \texttt{horizon=3} den
gewünschten $h$-Wert und \texttt{fields} die gewünschten Felder von $\rho(id,s)$
spezifiziert -- eine Verschmelzung der „deklarativen Datenanfrage“ von GraphQL mit der
„deklarativen Navigationsanfrage“ von HST, die formal als Erweiterung von
Definition~\ref{def:hst-repraesentation} um einen client-spezifizierten Parameter $h$
(statt eines fest konfigurierten) zu fassen wäre.

\subsubsection{Empirische Validierung des Kostenmodells $(\alpha,\beta,p)$}

Der in Abschnitt~\ref{sec:ausblick} (Unterabschnitt „Experimentelle Validierung in
Produktionssystemen“) für SC-REST skizzierte empirische Validierungsbedarf gilt für HST
in verstärktem Maße, da Satz~\ref{satz:optimaler-horizont} -- anders als die
asymptotischen Aussagen zu SC-REST -- von \emph{konkreten Werten} der Parameter
$\alpha,\beta,p,d$ abhängt. Eine produktive Instrumentierung müsste (1) $\alpha$ aus der
tatsächlich gemessenen Serialisierungsgröße von Knoten/Kanten kalibrieren (die in
Abbildung~\ref{fig:horizont-groesse} angenommenen 250\,Byte sind eine grobe Schätzung;
reale JSON- oder HAL-Repräsentationen könnten durch Kompression wie gzip oder durch das
in Abschnitt~\ref{sec:ausblick} (Graphkompression) skizzierte \emph{structure sharing}
deutlich darunter liegen), (2) $\beta$ aus der tatsächlichen Verteilung der
Round-Trip-Zeiten verschiedener Clientklassen (mobil, Desktop, Service-zu-Service)
erheben, und (3) $p$ aus den \emph{tatsächlichen} Sequenzlängenverteilungen
unterschiedlicher Prozesstypen schätzen -- letzteres ist methodisch eng mit der für
Abbildung~\ref{fig:lernkurve} verwendeten empirischen Fehlerratenkurve $p(n)$ verwandt,
jedoch mit umgekehrter Fragestellung: Während $p(n)$ die Wahrscheinlichkeit \emph{eines
Sequenzfehlers nach $n$ Trainingsanfragen} misst, beschreibt $p$ in
Definition~\ref{def:kostenmodell-horizont} die Wahrscheinlichkeit \emph{einer
Fortsetzung} der Sequenz unabhängig von deren Korrektheit. Mit kalibrierten
$(\hat\alpha,\hat\beta,\hat p)$ ließe sich die in
Abbildung~\ref{fig:optimaler-horizont} gezeigte theoretische Stufenfunktion $h^\ast$
gegen den empirisch tatsächlich kostenoptimalen Horizont validieren.

\subsubsection{Versionsmigration laufender Prozessinstanzen}

Satz~\ref{satz:hst-zustandslosigkeit} setzt voraus, dass $P$ während der Bearbeitung
einer Anfrage statisch ist; Abschnitt~\ref{sec:hst-cache} behandelt \emph{Erkennung} von
Änderungen $P\to P'$ via Signaturvergleich, jedoch nicht deren \emph{Behandlung} für
\emph{bereits laufende} Prozessinstanzen $(id,s)$ mit $s\in S$. Wird $P$ zu $P'$
aktualisiert (etwa weil ein neuer optionaler Teilschritt $\sigma_{r+1}$ eingeführt oder
ein Zustand $s\in S$ aus $S'$ entfernt wird), stellt sich die Frage, wie eine Instanz mit
$\pi(id)=s\notin S'$ zu behandeln ist. Eine formale Behandlung dieses Problems könnte
eine \emph{Migrationsabbildung} $\mu: S\rightharpoonup S'$ definieren, mit den
Anforderungen: (a) $\mu(s)=s$ für $s\in S\cap S'$ (unveränderte Zustände bleiben
erhalten), (b) für $s\in S\setminus S'$ muss $\mu(s)\in S'$ einen „äquivalenten“ Zustand
in $P'$ bezeichnen, wobei Äquivalenz bedeutet, dass die für $s$ in $P$ konformen
Restsequenzen (gemäß Definition~\ref{def:konformitaet}) unter $\mu$ auf in $P'$ ab
$\mu(s)$ konforme Restsequenzen abgebildet werden -- eine Bedingung, die strukturell an
die in Satz~\ref{satz:lokale-vollstaendigkeit} verwendete Charakterisierung über
$\delta^{*}$ erinnert, jedoch \emph{zwischen zwei verschiedenen Automaten} $P$ und $P'$
statt innerhalb eines einzelnen $G_s^h$. Existiert kein solches $\mu(s)$ (etwa weil ein
ganzer Prozesszweig in $P'$ entfernt wurde), müsste die Instanz in einen speziellen
Migrationszustand überführt und gegebenenfalls manuell oder durch eine
Kompensationslogik (vgl.\ das in Abschnitt~\ref{sec:ausblick} erwähnte Saga-Pattern und
\cite{garcia1987}) abgeschlossen werden. Die in Abschnitt~\ref{sec:hst-cache}
beschriebene signaturbasierte Erkennung von $P\to P'$ wäre damit der \emph{Trigger}, die
hier skizzierte Migrationsabbildung $\mu$ die \emph{Behandlung} -- zusammen ergäben sie
ein vollständiges Lebenszyklusmodell für Prozessautomaten, die sich über die Lebensdauer
einzelner Instanzen hinweg ändern können.

\subsection{Implementierungshinweis}

Der vollständige Code der Referenzimplementierung sowie ergänzende Testfälle sind
verfügbar unter: \url{https://github.com/hjstephan86/science/tree/main/science/screst}.

\begin{thebibliography}{99}

\bibitem{fielding2000}
Fielding, R.~T. (2000).
\emph{Architectural Styles and the Design of Network-based Software Architectures}.
Dissertation, University of California, Irvine.

\bibitem{rfc7231}
Fielding, R., Reschke, J. (Hrsg.) (2014).
\emph{Hypertext Transfer Protocol (HTTP/1.1): Semantics and Content}.
RFC~7231, IETF.

\bibitem{rfc7807}
Nottingham, M., Wilde, E. (2016).
\emph{Problem Details for HTTP APIs}.
RFC~7807, IETF.

\bibitem{richardson2007}
Richardson, L., Ruby, S. (2007).
\emph{RESTful Web Services}.
O'Reilly Media.

\bibitem{fowler2010}
Fowler, M. (2010).
\emph{Richardson Maturity Model -- steps toward the glory of REST}.
martinfowler.com.

\bibitem{webber2010}
Webber, J., Parastatidis, S., Robinson, I. (2010).
\emph{REST in Practice: Hypermedia and Systems Architecture}.
O'Reilly Media.

\bibitem{newman2021}
Newman, S. (2021).
\emph{Building Microservices: Designing Fine-Grained Systems}, 2.~Auflage.
O'Reilly Media.

\bibitem{hohpe2003}
Hohpe, G., Woolf, B. (2003).
\emph{Enterprise Integration Patterns: Designing, Building, and Deploying Messaging
Solutions}.
Addison-Wesley.

\bibitem{garcia1987}
Garcia-Molina, H., Salem, K. (1987).
\emph{Sagas}.
ACM SIGMOD Record, 16(3), S.~249--259.

\bibitem{helland2007}
Helland, P. (2007).
\emph{Life Beyond Distributed Transactions: An Apostate's Opinion}.
3rd Biennial Conference on Innovative Data Systems Research (CIDR).

\bibitem{abboud2015}
Abboud, A., Backurs, A., Vassilevska Williams, V. (2015).
\emph{Tight Hardness Results for LCS and Other Sequence Similarity Measures}.
IEEE 56th Annual Symposium on Foundations of Computer Science (FOCS).

\bibitem{epp2026subgraph}
Epp, S. (2026).
\emph{Der Subgraph Algorithmus}.
GitHub: \url{https://github.com/hjstephan86/subgraph}.

\bibitem{murata1989}
Murata, T. (1989).
\emph{Petri Nets: Properties, Analysis and Applications}.
Proceedings of the IEEE, 77(4), S.~541--580.

\bibitem{pautasso2008}
Pautasso, C., Zimmermann, O., Leymann, F. (2008).
\emph{RESTful Web Services vs. ``Big'' Web Services: Making the Right Architectural
Decision}.
Proceedings of the 17th International Conference on World Wide Web (WWW).

\end{thebibliography}

\end{document}

